\documentclass[JEP,XML,SOM,Unicode]{cedram}
\datereceived{2018-09-12}
\dateaccepted{2019-05-08}
\dateepreuves{2019-05-20}

\RequirePackage[mathscr]{eucal}
\usepackage{mathrsfs}
\let\mathcal \EuScript
\multlinegap0pt
\renewcommand{\theenumi}{\roman{enumi}}
\newcommand\mto{\mathchoice{\longmapsto}{\mapsto}{\mapsto}{\mapsto}}

\newcommand{\Psfrac}[2]{(\sfrac{#1}{#2})}
\newcommand{\psfrac}[2]{\sfrac{(#1)}{#2}}
\newcommand{\spfrac}[2]{\sfrac{#1}{(#2)}}

\newcommand{\scrL}{\mathscr{L}}
\newcommand{\scrO}{\mathscr{O}}
\newcommand{\scrQ}{\mathscr{Q}}

\newcommand{\cB}{\mathcal{B}}
\newcommand{\cC}{\mathcal{C}}
\newcommand{\cL}{\mathcal{L}}
\newcommand{\cQ}{\mathcal{Q}}
\newcommand{\cV}{\mathcal{V}}

\DeclareMathOperator{\dist}{dist}
\DeclareMathOperator{\Id}{Id}\let\sfId\Id
\DeclareMathOperator{\Span}{span}
\DeclareMathOperator{\Tr}{\mathsf{Tr}}

\newcommand{\intt}{\mathrm{int}}
\newcommand{\loc}{\mathrm{loc}}

\let\sep\mid

\providecommand{\enVert}[2][-1]{
\ensuremath{\mathinner{
\ifthenelse{\equal{#1}{-1}}{ % if
\!\left\lVert#2\right\rVert}{}
\ifthenelse{\equal{#1}{0}}{ % if
\lVert#2\rVert}{}
\ifthenelse{\equal{#1}{1}}{ % if
\!\bigl\lVert#2\bigr\rVert}{}
\ifthenelse{\equal{#1}{2}}{ % if
\!\Bigl\lVert#2\Bigr\rVert}{}
\ifthenelse{\equal{#1}{3}}{ % if
\!\biggl\lVert#2\biggr\rVert}{}
\ifthenelse{\equal{#1}{4}}{ % if
\!\Biggl\lVert#2\Biggr\rVert}{}
}}
}
\let\norm=\enVert

\usepackage{csquotes,dsfont,braket}

\theoremstyle{plain}
\newtheorem{theorem}{Theorem}[section]
\newtheorem{lemma}[theorem]{Lemma}
\newtheorem{proposition}[theorem]{Proposition}
\newtheorem{corollary}[theorem]{Corollary}
\theoremstyle{definition}
\newtheorem{notation}[theorem]{Notation}
\newtheorem{definition}[theorem]{Definition}
\newtheorem{remark}[theorem]{Remark}

\newcommand{\RR}{\mathbb{R}}
\newcommand{\CC}{\mathbb{C}}
\newcommand{\NN}{\mathbb{N}}
\newcommand{\dx}{\mathrm{\,d}}
\newcommand{\bn}{{\boldsymbol n}}
\newcommand{\bx}{{\boldsymbol x}}
\let\pa\partial
\let\eps\varepsilon
\DeclareMathOperator{\IM}{Im}\let\Im\IM
\DeclareMathOperator{\RE}{Re}\let\Re\RE
\newcommand{\Mat}{\mathsf{Mat}}
\newcommand{\jac}{\mathsf{Jac}}
\DeclareMathOperator{\supp}{\mathsf{supp}}
\DeclareMathOperator{\dive}{div}
\DeclareMathOperator{\Dom}{\mathsf{Dom}}

\begin{document}


\title[An infinite mass limit]{The MIT Bag Model as an infinite mass limit}

\author[\initial{N.} \lastname{Arrizabalaga}]{\firstname{Naiara} \lastname{Arrizabalaga}}
\address{Departamento de Matemáticas, Universidad del País Vasco/Euskal Herriko Unibertsitatea (UPV/EHU)\\
48080 Bilbao, Spain}
\email{naiara.arrizabalaga@ehu.eus}

\author[\initial{L.} \lastname{Le Treust}]{\firstname{Loïc} \lastname{Le Treust}}
\address{Aix Marseille Univ, CNRS, Centrale Marseille, I2M\\
Marseille, France}
\email{loic.le-treust@univ-amu.fr}
\urladdr{https://sites.google.com/site/loicletreust/}

\author[\initial{A.} \lastname{Mas}]{\firstname{Albert} \lastname{Mas}}
\address{Departament de Matemàtiques, Universitat Politècnica de Catalunya\\
Campus Diagonal Besòs, Edifici A (EEBE), Av. Eduard Maristany 16, 08019 Barcelona, Spain}
\email{albert.mas.blesa@upc.edu}
\urladdr{https://mat-web.upc.edu/people/albert.mas.blesa/}

\author[\initial{N.} \lastname{Raymond}]{\firstname{Nicolas} \lastname{Raymond}}
\address{Laboratoire Angevin de Recherche en Mathématiques, LAREMA, UMR 6093, UNIV Angers, SFR Math-STIC\\
2, boulevard Lavoisier, 49045 Angers Cedex 01, France}
\email{nicolas.raymond@univ-angers.fr}
\urladdr{https://nraymond.perso.math.cnrs.fr/}

\thanks{N.A.\ was partially supported by ERCEA Advanced Grant 669689-HADE, MTM2014-53145-P (MICINN, Gobierno de España) and IT641-13 (DEUI, Gobierno Vasco).
L.LT.\ was partially supported by ANR DYRAQ ANR-17-CE40-0016-01.
A.M.\ was partially supported by
MTM2017-84214 and MTM2017-83499 projects of the MCINN (Spain),
2017-SGR-358 project of the AGAUR (Catalunya), and ERC-2014-ADG project
HADE Id.\! 669689 (European Research Council).}

\keywords{Dirac operator, relativistic particle in a box, MIT bag model, spectral theory}

\subjclass{35J60, 35Q75, 49J45, 49S05, 81Q10, 81V05, 35P15, 58C40}

\begin{abstract}
The Dirac operator, acting in three dimensions, is considered. Assuming that a large mass $m>0$ lies outside a smooth enough and bounded open set $\Omega\subset\mathbb{R}^3$, it is proved that its spectrum approximates the one of the Dirac operator on $\Omega$ with the MIT bag boundary condition. The approximation, modulo an error of order $o(1/\sqrt m)$, is carried out by introducing tubular coordinates in a neighborhood of $\partial\Omega$ and analyzing one dimensional optimization problems in the normal direction.
\end{abstract}

\alttitle{Le modèle MIT bag obtenu comme une limite de masse grande}

\altkeywords{Opérateur de Dirac, particules relativistes dans une boîte, modèle MIT bag, théorie spectrale}

\begin{altabstract}
Nous considérons l'opérateur de Dirac en dimension $3$ dont la masse $m>0$ est supposée grande à l'extérieur d'un ouvert borné et régulier $\Omega\subset\mathbb{R}^3$. Nous démontrons que son spectre approche celui de l'opérateur de Dirac sur $\Omega$ qui intègre dans son domaine les conditions au bord dites \og MIT bag\fg. L'analyse asymptotique est réalisée grâce à l'usage de coordonnées tubulaires et à l'analyse d'un problème d'optimisation unidimensionnel dans la direction normale.
\end{altabstract}

\maketitle
\tableofcontents


\section{Introduction}
\subsection{Context}
This paper is devoted to the spectral analysis of the Dirac operator with high scalar potential barrier in three dimensions. More precisely, we will assume that there is a large mass $m$ outside a smooth and bounded open set $\Omega$. From physical considerations, see \cite{bogoliubov1987,MIT061974}, it is expected that, when $m$ becomes large, the eigenfunctions of low energy do not visit $\RR^3\setminus\Omega$ and tend to satisfy the so-called MIT bag condition on $\partial\Omega$. This boundary condition, that we will define in the next section, is usually chosen by the physicists \cite{johnson,MIT061974,PhysRevD.12.2060}, in order to get a vanishing normal flux at the bag surface. It was originally introduced by Bogolioubov in the late $60's$ \cite{bogoliubov1987} to describe the confinement of the quarks in the hadrons with the help of an infinite scalar potential barrier outside a fixed set $\Omega$. In the mid $70's$, this model has been revisited into a shape optimization problem named \emph{MIT bag model} \cite{johnson,MIT061974,PhysRevD.12.2060} in which the optimized energy takes the form
\[
\Omega\mto \lambda_1(\Omega) + b|\Omega|,
\]
where $\lambda_1(\Omega)$ is the first nonnegative eigenvalue of the Dirac operator with the boundary condition introduced by Bogolioubov, $|\Omega|$ is the volume of $\Omega\subset\RR^3$ and $b>0$.
The interest of the bidimensional equivalent of this model has recently been renewed with the study of graphene where this condition is sometimes called \enquote{infinite mass condition}, see \cite{PhysRevB.77.085423,MR901725}. The aim of this paper is to provide a mathematical justification of this terminology, and extend to dimension three the work \cite{stockmeyer2016infinite}.
More precisely, we show the convergence of the eigenvalues for the Dirac operator with high scalar potential barrier to the ones of the MIT bag Dirac operator. In dimension two, this follows by the convergence of the spectral projections shown in \cite{stockmeyer2016infinite}. Regarding the first eigenvalue of the MIT bag Dirac operator, we also find the first order term in the asymptotic expansion of the eigenvalues given by the high scalar potential barrier, showing its dependence on geometric quantities related to $\partial\Omega$. This is a novel result with respect to the ones in \cite{stockmeyer2016infinite}.

\subsection{The Dirac operator with large effective mass}

In the whole paper, $\Omega$ deno\-tes a fixed bounded domain of $\RR^3$ with $\cC^{2,1}$ boundary.

Let us recall the definition of the Dirac operator associated with the energy of a relativistic particle of mass $m_0\in \RR$ and spin $\sfrac{1}{2}$, see \cite{Thaller1992}. The Dirac operator is a first order differential operator $(H,\Dom(H))$, acting on $L^2(\RR^3;\mathbb{C}^4)$ in the sense of distributions, defined by
\begin{equation}
H=c\alpha\cdot D+m_0c^2\beta,\qquad D=-i\hbar\nabla,
\end{equation}
where $\Dom(H) = H^1(\RR^3; \CC^4)$, $c>0$ is the velocity of light, $\hbar>0$ is Planck's constant, $\alpha=(\alpha_1,\alpha_2,\alpha_3)$ and $\beta$
are the $4\times4$ Hermitian and unitary matrices given by
\[
\beta=\begin{pmatrix}1_2&0\\0&-1_2\end{pmatrix},\quad
\alpha_k=\begin{pmatrix}0&\sigma_k\\\sigma_k&0\end{pmatrix}\quad\mbox{for }k=1,2,3.
\]

Here, the Pauli matrices $\sigma_1,\sigma_2$ and $\sigma_3$ are defined by
\[\arraycolsep3.5pt
\sigma_1 = \begin{pmatrix}
0&1\\1&0
\end{pmatrix},
\quad
\sigma_2 = \begin{pmatrix}0&-i\\i&0
\end{pmatrix},
\quad
\sigma_3 = \begin{pmatrix}
1&0\\0&-1
\end{pmatrix},
\]
and $\alpha \cdot X$ denotes $\sum_{j=1}^3\alpha_j X_j$ for any $X = (X_1,X_2,X_3)$.
In the following, we shall always use units with $\hbar=c=1$.

The Dirac and Pauli matrices are chosen in such a way that the Dirac operator $(H,\Dom(H))$ is self-adjoint, and satisfies
\[
H^2 = 1_4(m_0^2-\Delta),
\]
(see for instance \cite[\S 1.1]{Thaller1992}).
Let us also mention that its spectrum is
\[(-\infty,-|m_0|]\cup[|m_0|,+\infty).\]

In this paper, we consider particles with large effective mass $m\gg m_0$ outside~$\Omega$. Their kinetic energy is associated with the self-adjoint operator $(H_m, \Dom(H_m))$ \hbox{defined} by
\[
H_{m}=\alpha\cdot D+(m_0 + m\chi_{\Omega'})\beta,
\]
where $\Omega'$ is the complementary set of $\overline{\Omega}$, $\chi_{\Omega'}$ is the characteristic function of $\Omega'$ and $\Dom(H_m) = H^1(\RR^3;\CC^4)$.
The essential spectrum of $(H_m, \Dom(H_m))$ is
\[
(-\infty,-|m_0 + m|]\cup[|m_0 + m|,+\infty).
\]
In this paper, the mass $m_0$ is not assumed to be positive since this assumption is not used in the proofs (see also Remark \ref{rem.opplimit}).

\begin{notation}\label{not.def1}
In the following, $\Gamma := \pa \Omega$ and for all $\bx\in\Gamma$, $\bn(\bx)$ is the outward-pointing unit normal vector to the boundary, $L(\bx)=d\bn_{\bx}$ denotes the second fundamental form of the boundary, and
\begin{equation*}
\kappa(\bx)=\Tr L(\bx) \quad\text{and}\quad K(\bx) = \det L(\bx)
\end{equation*}
are the mean curvature and the Gauss curvature of $\Gamma$, respectively.
\end{notation}

\begin{definition}\label{def.MIT}
The MIT bag Dirac operator $(H^\Omega,\Dom(H^\Omega))$ is defined on the domain
\[
\Dom(H^\Omega) = \{\psi\in H^1(\Omega;\mathbb{C}^4)\sep\cB\psi = \psi~\text{on}~\Gamma\},\qquad\mbox{with}\quad\cB=-i\beta(\alpha\cdot \bn),
\]
by $H^\Omega\psi=H\psi$ for all $\psi\in\Dom(H^\Omega)$. Observe that the trace is well-defined by a classical trace theorem.

If $\Gamma$ is $\cC^2$, the operator $(H^\Omega,\Dom(H^\Omega))$ is self-adjoint with compact resolvent \cite{OBV17,ALR17SA,booss2009calderon,benguria2016spectral,BB12}.
\end{definition}
\begin{notation}
We denote by $\braket{\cdot\,,\cdot}$ the $\CC^4$ scalar product (antilinear with respect to the left argument) and by $\braket{\cdot\,,\cdot}_{U}$ the $L^2$ scalar product on the set $U\subset\RR^3$.
\end{notation}

\begin{notation}
We define, for every $\bn\in \mathbb{S}^2$, the orthogonal projections
\begin{equation}\label{eq:boundary_projections}
\Xi^\pm = \frac{1_4\pm\cB}{2}
\end{equation}
associated with the eigenvalues $\pm1$ of the matrix $\cB$.
\end{notation}

\subsection{Squared operators, heuristics, and main results}
The aim of this paper is to relate the spectra of $H_{m}$ and $H^\Omega$ in the limit $m\to+\infty$.

\begin{notation}\label{not.lambda}

Let $(\lambda_k)_{k\in \NN^*}$ and $(\lambda_{k,m})_{k\in \NN^*}$ be
the increasing sequences defined~by
\begin{equation*}
\begin{split}
\lambda_k &= \inf_{\substack{
V\subset \Dom(H^\Omega)
\\
\dim V = k}}\;
\sup_{\substack{
\varphi\in V
\\
\norm{ \varphi}_{L^2(\Omega)} = 1}}
\norm{H^\Omega \varphi}_{L^2(\Omega)}
\\
&=
\sup_{\{\psi_{1},\ldots, \psi_{k-1}\}\subset \Dom(H^\Omega)}\;
\inf_{\substack{\varphi\in \Span(\psi_{1},\ldots, \psi_{k-1})^\perp
\\
\norm{ \varphi}_{L^2(\Omega)} = 1}}
\norm{H^\Omega \varphi}_{L^2(\Omega)},
\end{split}
\end{equation*}
and
\begin{equation*}
\begin{split}
\lambda_{k,m} &= \inf_{\substack{V\subset H^1(\RR^3;\CC^4)
\\
\dim V = k}}\;
\sup_{\substack{\varphi\in V
\\
\norm{ \varphi}_{L^2(\RR^3)} = 1}}
\norm{H_m \varphi}_{L^2(\RR^3)}
\\
&=
\sup_{\{\psi_{1},\ldots, \psi_{k-1}\}\subset H^1(\RR^3;\CC^4)}\;
\inf_{\substack{\varphi\in \Span(\psi_{1},\ldots, \psi_{k-1})^\perp
\\
\norm{ \varphi}_{L^2(\RR^3)} = 1}}
\norm{H_m \varphi}_{L^2(\RR^3)},
\end{split}
\end{equation*}
for $k\in \NN^*$ and $m>0$. Here, $\NN^* := \NN\setminus\{0\}$. By the $\min$-$\max$ characterization and the properties given in Definition \ref{def.MIT}, the sequence $(\lambda_k)_{k\in \NN^*}$ is made of all the eigenvalues of the operator $|H^\Omega|$, each one being repeated according to its multiplicity. Similarly, the terms of the sequence $(\lambda_{k,m})_{k\in \NN^*}$ that satisfy
\[
\lambda_{k,m}<|m_0+m|\,
\]
are the eigenvalues of $|H_m|$ lying below its essential spectrum $[|m_0+m|,+\infty)$, each one being repeated according to its multiplicity. For $k$ large enough, this sequence may become stationary at $|m_0+m|$.
\end{notation}

\subsubsection{The quadratic forms}
At first sight, it might seem surprising that $\lambda_{k}$ and $\lambda_{k,m}$ are related, especially because of the boundary condition of $H^\Omega$. It becomes less surprising when computing the squares of the operators. This is the purpose of the following lemma.

\begin{lemma}\label{lem:quadf1}
Let $\varphi \in \Dom(H^\Omega)$ and $\psi\in H^1(\RR^3; \CC^4)$. Then
\begin{equation}\label{eq:quadMIT}
\|H^\Omega \varphi\|^2_{L^2(\Omega)} = \cQ^{\intt}(\varphi) := \|\nabla \varphi\|^2_{L^2(\Omega)} + \int_{\Gamma}\left(\frac{\kappa}{2}+m_0\right)|\varphi|^2\dx \Gamma + m_0^2\norm{\varphi}_{L^2(\Omega)}^2,
\end{equation}
where $\kappa$ is defined in Notation \ref{not.def1},
and
\begin{equation}\label{eq:quadIM}\begin{split}
\|H_m \psi\|^2_{L^2(\RR^3)} &=
\|\nabla\psi\|_{L^2(\Omega)}^2
+\|\nabla\psi\|_{L^2(\Omega')}^2 + \|(m_0 +m\chi_{\Omega'}) \psi\|_{L^2(\RR^3)}^2\\
&\hspace*{4.3cm}-m\Re\langle\cB\psi, \psi\rangle_{\Gamma}\\
&=\|\nabla\psi\|_{L^2(\Omega)}^2
+\|\nabla\psi\|_{L^2(\Omega')}^2 + \|(m_0 +m\chi_{\Omega'}) \psi\|_{L^2(\RR^3)}^2\\
&\hspace*{4.3cm}+m\|\Xi^-\psi\|^2_{L^2(\Gamma)}-m\|\Xi^+\psi\|^2_{L^2(\Gamma)}.
\end{split}
\end{equation}
\end{lemma}
\begin{proof}
The equality \eqref{eq:quadMIT} is proved for instance in \cite[\S A.2]{MR3663620}. Let $\psi\in H^1(\RR^3; \CC^4)$. By integrations by parts,
\[
\begin{split}
\|H_m\psi\|_{L^2(\RR^3)}^2 &
=
\|\alpha\cdot D\psi\|_{L^2(\RR^3)}^2 + \|(m_0 +m\chi_{\Omega'}) \psi\|_{L^2(\RR^3)}^2
+2m\Re\langle\alpha\cdot D \psi, \beta \psi\rangle_{\Omega'}
\\
&=
\|\nabla\psi\|_{L^2(\RR^3)}^2 + \|(m_0 +m\chi_{\Omega'}) \psi\|_{L^2(\RR^3)}^2
-m\Re\langle\cB\psi, \psi\rangle_{\Gamma}.
\end{split}
\]
Then, note that, for all $\psi\in H^1(\RR^3;\CC^4)$,
\[\Re\langle\cB\psi, \psi\rangle_{\Gamma}=\|\Xi^+\psi\|^2_{L^2(\Gamma)}-\|\Xi^-\psi\|^2_{L^2(\Gamma)}.\qedhere\]
\end{proof}

Considering \eqref{eq:quadIM} leads to the following minimization problem, for $v\in H^{1}(\Omega)$,
\begin{equation}\label{pb:extopt}
\Lambda_m(v) = \inf \{
\cQ_m(u), u\in V_v
\},\quad \cQ_m(u) = \|\nabla u\|^2_{L^2(\Omega')} + m^2\|u\|^2_{L^2(\Omega')},
\end{equation}
where
\[V_v = \{u\in H^1(\Omega',\mathbb{C}^4) \sep u = v\mbox{ on }\Gamma\}.\]

A classical extension theorem (see \cite[\S 5.4]{evans1998partial}) ensures that $V_v$ is non-empty.
\subsubsection{Heuristics}\label{sec:heuri}
In this paper, we will analyze the behavior of $\Lambda_{m}(v)$ and prove in particular (see Proposition \ref{prop:ext}) that there exists $C>0$ such that for $m$ large, and all $v\in H^1(\Omega;\CC^4)$
\begin{equation}\label{eq.Lambdam1}
o(1)\geq\Lambda_{m}(v)-\left(m\|v\|^2_{L^2(\Gamma)} + \int_\Gamma \frac{\kappa}{2}\,|v|^2\dx \Gamma\right)\geq
-\frac{C}{m}\,\|v\|^2_{H^1(\Gamma)}.
\end{equation}
Replacing $m$ by $m_0+m$ in \eqref{eq.Lambdam1}, we get, for all $\psi\in H^1(\RR^3;\CC^4)$,
\begin{multline}\label{eq.lbHm}
\|H_m\psi\|_{L^2(\RR^3)}^2\geq \|\nabla\psi\|^2_{L^2(\Omega)} + m_0^2\norm{\psi}_{L^2(\Omega)}^2
\\
+ \int_\Gamma\left(\frac{\kappa}{2}+m_0\right)|\psi|^2\dx \Gamma+2m\|\Xi^-\psi\|^2_{L^2(\Gamma)} - \frac{C}{m}\,\|\psi\|^2_{L^2(\Gamma)}.
\end{multline}
Take any eigenfunction $\varphi$ of $H^\Omega$ and consider a minimizer $u_{\varphi}$ of \eqref{pb:extopt} for $v=\varphi$ and~$m$ replaced by $m+m_0$.
Then, letting $\psi=\mathds{1}_{\Omega} \varphi+\mathds{1}_{\Omega'}u_{\varphi}\in H^1(\RR^3;\CC^4)$,
we get
\[\|H_{m}\psi\|_{L^2(\RR^3)}^2
=
\|\nabla\varphi\|^2_{L^2(\Omega)}
+ m_0^2\norm{\psi}_{L^2(\Omega)}^2
+\Lambda_{m+m_0}(\varphi)-m\|\Xi^+\varphi\|_{L^2(\Gamma)}^2.\]
With \eqref{eq.Lambdam1} at hand, we deduce that, for all $j\in\NN^*$,
\[\lambda^2_{j,m}\leq \lambda^2_{j}+o(1).\]

Conversely, if we are interested in the eigenvalues of $(H_{m})^2$ that are of order $1$ when $m\to+\infty$, we see from \eqref{eq.lbHm} that the corresponding normalized eigenfunctions must satisfy $\Xi^-\psi=\scrO(m^{-1})$ and, in particular, $\cB\psi=\psi+\scrO(m^{-1})$. Thus, we get formally, for all $j\in\NN^*$,
\[\lambda^2_{j,m}\geq \lambda^2_{j}+o(1).\]
The aim of this paper is to make this heuristics rigorous. We now state our main theorem.

\begin{theorem}\label{thm:main}
Let $\Omega\subset\RR^3$ be a bounded domain of class $\cC^{2,1}$ (\ie the derivative of the curvatures is bounded).
The singular values of $H_{m}$ can be estimated as follows:
\begin{enumerate}
\item $\lim_{m\to+\infty} \lambda_{k,m} = \lambda_k$, for all $k\in \NN^*$.

\item Let $k_1\in \NN^*$ be the multiplicity of the first eigenvalue $\lambda_1$ of $|H^\Omega|$. For all $k\in\{1,\dots,k_1\}$, we have
\[
\lambda_{k,m} = \left(\lambda_1^2+\frac{\nu_{k}}{m}+o\left(\sfrac{1}{m}\right)\right)^{1/2},
\]
where
\begin{equation}\label{eq:varmu}
\nu_{k} = \inf_{\substack{V\subset \ker(|H^\Omega|-\lambda_1)\\
\dim V = k}}\;
\sup_{\substack{u\in V\\
\norm{u}_{L^2(\Omega)}=1}}
\eta(u),
\end{equation}
with
\[
\eta(u) =
\int_\Gamma \Bigl(
\frac{|\nabla_s u|^2}{2}
-
\frac{|(\pa_n+\kappa/2+m_0)u|^2}{2}
+
\Bigl(\frac{K}{2}-\frac{\kappa^2}{8}-\frac{\lambda_1^2}{2}\Bigr)|u|^2
\Bigr)\dx\Gamma.
\]
\end{enumerate}
Here, $(\lambda_k)_{k\in \NN^*}$ and $(\lambda_{k,m})_{k\in \NN^*}$ are defined in Notation \ref{not.lambda}, and
$\kappa$ and $K$ are defined in Notation \ref{not.def1}. We denote by $\pa_\bn$ the outward pointing normal derivative and by $\nabla_s$ the tangential gradient on $\Gamma$.
\end{theorem}
\begin{remark}
The max-min formula \eqref{eq:varmu} makes sense since $\ker(|H^\Omega|-\lambda \Id)\subset H^2(\Omega;\CC^4)$ for any eigenvalue $\lambda$ of $|H^\Omega|$.
\end{remark}
\begin{remark}
$H_m$ and $H^\Omega$ anticommute with the charge conjugation $C$ defined, for all $\psi\in \CC^4$, by
\[
C\psi = i\beta \alpha_2 \overline{\psi},
\]
where $\overline{\psi}\in \CC^4$ is the vector obtained after complex conjugations of each of the components of $\psi$ (see for instance \cite[\S 1.4.6]{Thaller1992} and \cite[\S A.1]{MR3663620}). As a consequence, the spectrum of $H_m$ and $H^\Omega$ are symmetric with respect to $0$, and Theorem \ref{thm:main} may be rewritten as a result on the eigenvalues of $H_m$ and $H^\Omega$.
\end{remark}

\begin{remark}\label{rem.opplimit}
Let us define the operator $(\widetilde{H^\Omega},\Dom(\widetilde{H^\Omega}))$ on
\[
\Dom(\widetilde{H^\Omega}) = \{\psi\in H^1(\Omega;\mathbb{C}^4)\sep\cB\psi = -\psi~\text{on}~\Gamma\}\,
\]
by $\widetilde{H^\Omega}\psi=H\psi$ for all $\psi\in\Dom(\widetilde{H^\Omega})$. It is the MIT bag Dirac operator with \emph{reversed} boundary condition (see Definition \ref{def.MIT}). The singular values of $\widetilde{H^\Omega}$ are approximated by the singular values of $H_m$ as $m$ tends to $-\infty$. This follows immediately from Theorem \ref{thm:main}, conjugating all the operators by the chirality matrix
\[\arraycolsep3.5pt
\gamma_5=\begin{pmatrix}0&1_2\\1_2&0\end{pmatrix},
\]
and by using the algebraic properties
\begin{equation*}
\begin{split}
\beta \gamma_5 = -\gamma_5\beta,\quad
\gamma_5(\alpha\cdot \bx) = (\alpha\cdot \bx)\gamma_5,\quad
\gamma_5\cB\gamma_5 = -\cB,
\end{split}
\end{equation*}
for all $\bx\in\RR^3$.
\end{remark}
\begin{remark}
Our proof of Theorem \ref{thm:main} also provides the convergence of the eigenprojectors associated with the first eigenvalues of $|H_m|$. They converge towards the eigenprojectors associated with the first eigenvalues of $|H^\Omega|$, see Lemma \ref{lem:induc2} and Remark \ref{rem.convergencefunctions}, and \cite[Th.\,1]{stockmeyer2016infinite} for the two-dimensional case.
\end{remark}
\begin{remark}
In view of Theorem~\ref{thm:main}, it is natural to ask if one has convergence of $H_m$ to $H^\Omega$ in some resolvent sense when $m\to+\infty$. On one hand, in the recent work \cite{barbaroux2018normconvergence} it is shown the convergence in the norm resolvent sense for the bidimensional analogues of $H_m$ and $H^\Omega$. On the other hand, in \cite{maspizzichillo2018} the authors study interactions of the free Dirac operator in $\RR^3$ with potentials that shrink towards $\partial\Omega$, proving convergence in the strong resolvent sense to $\delta$-shell interactions with precise coupling constants. As $m\to+\infty$, our operator $H_m$ may be seen as a degenerate case of the interactions with shrinking potentials considered in \cite{maspizzichillo2018} and, at a formal level, in this case the resulting $\delta$-shell interaction leads to the operator $H^\Omega$.

The above-mentioned results suggest that convergence in the norm (or at least strong) resolvent sense may also hold in our three dimensional setting.
\end{remark}

\subsubsection{A vectorial Laplacian with Robin-type boundary conditions}\label{sec.robintro}
Let us also mention an intermediate spectral problem whose study is needed in our proof of Theorem~\ref{thm:main} and that may be of interest on its own.
We consider the vectorial Laplacian associated with the quadratic form
\begin{equation}\label{eq:quadf2}
\begin{split}
\cQ_m^\intt(u) = \norm{\nabla u}_{L^2(\Omega)}^2 +m_0^2\norm{u}_{L^2(\Omega)}^2+ \int_\Gamma \left(\frac{\kappa}{2}+m_0\right)|u|^2\dx \Gamma + 2m\norm{\Xi^- u}_{L^2(\Gamma)}^2
\end{split}
\end{equation}
for $u\in \Dom(\cQ_m^\intt) = H^1(\Omega; \CC^4)$ and $m>0$, where $\Xi^-, \Xi^+$ are defined by \eqref{eq:boundary_projections}. By a classical trace theorem, this form is bounded from below. More precisely, we have the following result whose proof is sketched in Section \ref{sec.Lintdef}.
\begin{lemma}\label{lem:sarobint}
The self-adjoint operator associated with $\cQ_m^\intt$ is defined by
\begin{equation}\label{eq:laplacian_int}
\begin{split}
\Dom(L_m^\intt) &= \left\{u\in H^2(\Omega; \CC^4) \Bigm| \begin{array}{l}
\Xi^-\left(\pa_\bn + \kappa/2+m_0 + 2m\right) u =0 \mbox{ on }\Gamma,\\
\Xi^+\left(\pa_\bn + \kappa/2+m_0\right) u = 0 \mbox{ on }\Gamma
\end{array}\right\}
\\
L_m^\intt u &= \left(-\Delta+m_0^2\right) u \quad\mbox{for all } u \in \Dom(L_m^\intt).
\end{split}
\end{equation}
It has compact resolvent and its spectrum is discrete.
\end{lemma}

Using an integration by parts and the identities \eqref{eq:boundary_projections}, we get
\[
\braket{u,L_m^\intt u}_\Omega = \cQ_m^\intt(u),
\]
for all $u\in\Dom(L_m^\intt)$.

\begin{notation}\label{not.lambdaint}
Let $(\lambda_{k,m}^\intt)_{k\in \NN^*}$ denote the sequence of eigenvalues, each one being repeated according to its multiplicity and such that
\begin{equation}\label{eq:eigen_laplacian_int}
\lambda_{1,m}^\intt\leq \lambda_{2,m}^\intt\leq \cdots
\end{equation}
\end{notation}
The asymptotic behavior of the eigenvalues of $L_m^\intt$ is detailed in the following theorem.
\begin{theorem}\label{thm:asym_eigen_int-0}
The following holds:
\begin{enumerate}
\item For every $k\in \NN^*$, $\lim_{m\to +\infty} \lambda_{k,m}^\intt = \lambda_k^2$.

\item Let $\lambda$ be an eigenvalue of $|H^\Omega|$ of multiplicity $k_1\in \NN^*$. Consider the unique integer $k_0\in \NN$ such that for all $k\in\{1,\dots,k_1\}$, $\lambda_{k_0+k} = \lambda$.

Then, for all $k\in\{1,2,\dots, k_1\}$, we have
\[
\lambda_{k_0+k,m}^\intt = \lambda^2 +\frac{\mu_{\lambda,k}}{m} + o\left(\sfrac{1}{m}\right),
\]
where
\begin{equation}\label{eq:mulk}
\mu_{\lambda,k} :=
\inf_{\substack{V\subset \ker(|H^\Omega|-\lambda)\\
\dim V = k}}
\sup_{\substack{v\in V\\
\norm{v}_{L^2(\Omega)}=1}}
-\frac{\norm{(\pa_\bn+\kappa/2+m_0)v}_{L^2(\Gamma)}^2}{2}.
\end{equation}
\end{enumerate}
Here, $(\lambda_k)_{k\in \NN^*}$ is defined in Notation \ref{not.lambda}, $(\lambda_{k,m}^\intt)_{k\in \NN^*}$ in Notation \ref{not.lambdaint}, and $\kappa$ in Notation \ref{not.def1}.
\end{theorem}

\subsection{Organization of the paper}
In Section \ref{sec.2}, we discuss the asymptotic properties of the minimizers associated with the exterior optimization problem \eqref{pb:extopt}.
In Section \ref{sec.3}, we investigate the interior problem given by \eqref{eq:quadf2}. Finally, in Section \ref{sec.4}, we prove Theorem \ref{thm:main}.

In order to ease the reading, we provide here a list of notation regarding the spaces and the quadratic forms, as well as the equation number where they are introduced, that we will use in the sequel:
\[\arraycolsep3.5pt
\begin{array}{|l|c|c|c|c|}
\hline
\mbox{Key}&\mbox{Space domain} &\mbox{Variational space}& \mbox{Quadratic form}&\mbox{Infimum}\\
\hline
\mbox{\eqref{eq:quadMIT}}&\Omega &\Dom(H^\Omega)&\cQ^\intt&-\\
\mbox{\eqref{pb:extopt}}&\Omega' &V_v&\cQ_m&\Lambda_m(v)\\
\mbox{\eqref{eq:quadf2}}&\Omega &H^1(\Omega; \CC^4)&\cQ_m^\intt&-\\
\mbox{\eqref{pb:extopt1}}&\mathcal V_{\delta} &V_{v,\delta}& \cQ_{m}& \Lambda_{m,\delta}(v)\\
\mbox{\eqref{eq:dom-hat}}&\widehat{\mathcal V}_{m} &\widehat V_{m} &\widehat{\scrQ}_{m}&\Lambda_{m,m^{-1/2}}(v)\\
\mbox{\eqref{eq:optiextpar}}&(0,\sqrt{m}) & \widehat{V}_{m,\kappa,K} & \widehat{\scrQ}_{m,\kappa,K}& \Lambda_{m,\kappa,K} \\
\hline
\end{array}
\]

\section{About the exterior optimization problem}\label{sec.2}
The aim of this section is to study the minimizers of \eqref{pb:extopt} and their properties when $m$ tends to $+\infty$. These properties are gathered in the following proposition.

\begin{proposition}\label{prop:ext}
For all $v\in H^1(\Omega)$, there exists a unique minimizer $u_m(v)$ associated with $\Lambda_m(v)$, and it satisfies, for all $u\in V_v$,
\[\cQ_m(u) = \Lambda_m(v) + \cQ_m(u-u_m(v)).\]

Moreover, the following holds:
\begin{enumerate}
\item\label{item.i2}
Assume that $\Gamma$ is $\cC^2$. There exist $C,m_1>0$ such that, for every $m\geq m_1$, $v\in H^1(\Omega)$,
\[
Cm\norm{v}_{H^1(\Omega)}^2
\geq\Lambda_{m}(v)
\geq \left(m\|v\|^2_{L^2(\Gamma)} + \int_\Gamma \frac{\kappa}{2}\,|v|^2\dx \Gamma\right)
-\frac{C}{m}\,\|v\|^2_{L^2(\Gamma)}.
\]
\end{enumerate}
Assume that $\Gamma$ is $\cC^{2,1}$. There exists $C>0$ such that, for every $m\geq m_1$,
\begin{enumerate}\refstepcounter{enumi}
\item\label{item.i4}
for $v\in H^1(\Omega)$,
\[\begin{split}
\left(m\|v\|^2_{L^2(\Gamma)} + \int_\Gamma \frac{\kappa}{2}\,|v|^2\dx \Gamma\right)+o(1)&\geq\Lambda_{m}(v).
\end{split}\]
Here, the term $o(1)$ depends on $v$ (not only on the $H^1$ norm of $v$).
\item\label{item.ii} for all $v\in H^2(\Omega)$,
\[
\left|\Lambda_m(v)-\widetilde\Lambda_m(v)\right|\leq \frac{C}{m^{3/2}}\,\|v\|^2_{H^{3/2}(\Gamma)},
\]
\item\label{item.iii} for all $v\in H^2(\Omega)$,
\begin{gather*}
\Bigl|\|u_m(v)\|^2_{L^2(\Omega')} - \frac{\|v\|^2_{L^2(\Gamma)}}{2m}\Bigr|\leq \frac{C}{m^2}\,\|v\|^2_{H^{3/2}(\Gamma)},
\\
\widetilde \Lambda_{m}(v) = m\int_\Gamma
|v|^2
\dx\Gamma
+
\int_\Gamma
\frac{\kappa}{2}\,|v|^2
\dx\Gamma
+
m^{-1}\int_\Gamma \Big\{
\frac{|\nabla_s v|^2}{2}
+
\Big(\frac{K}{2}-\frac{\kappa^2}{8}\Big)|v|^2
\Big\}\dx\Gamma.
\end{gather*}
\end{enumerate}
\end{proposition}

\subsection{Organization of the section}
Since there are many steps in the proof of Proposition \ref{prop:ext}, let us briefly describe the strategy:

\begin{itemize}
\item In Section \ref{sec.exist-unique}, we explain why the minimizers exist, are unique, and we describe their Euler-Lagrange equations.
\item In Section \ref{sec.Agmon}, we prove Proposition \ref{lem:agmon1}. This proposition states that, when $m$ goes to $+\infty$, the minimizers are exponentially localized near the interface $\Gamma$. This allows to replace our optimization problem on $\Omega'$ by the same optimization problem on a thin (of size $m^{-1/2})$ neighborhood of $\Gamma$.
\item In Section \ref{sec.tub}, we study the optimization problem in the tubular neighborhood. In this \enquote{tube}, we can use the classical tubular coordinates, called $(s,t)$, where $s\in\Gamma$ and $t$ represents the distance to $\Gamma$. In these coordinates, we are led to consider a \enquote{transverse} optimization problem, that is a problem in one dimension (with respect to $t$) with parameters involving the curvature of the boundary. Then, explicit computations provide the asymptotics of the $1D$-minimizers.
\item In Section \ref{sec.proof-prop}, we establish Proposition \ref{prop:ext}. In particular, we use the projection on the $1D$-minimizers to give the asymptotics of the minimizers in the tubular neighborhood. Note that our refined bounds are proved under the assumption that the boundary is of class $\cC^{2,1}$. Indeed, we need at least $\cC^{2,1}$ regularity to control the tangential derivative of the transverse optimizers (which depend on the curvature, see Lemma \ref{lem.deriv-u}) when establishing, for instance, the accurate upper bound of $\Lambda_m(v)$ (see Corollary \ref{cor:estilowreg}).
\end{itemize}

\subsection{Existence, uniqueness and Euler-Lagrange equations}\label{sec.exist-unique}

Let us discuss here the existence of the minimizers announced in Proposition \ref{prop:ext} and their elementary properties. We will see later that, in the limit $m\to+\infty$, this minimization problem on $\Omega'$ is closely related to the same problem on a tubular neighborhood in $\Omega'$ of $\Gamma$. For $\delta>0$, $m>0$, and $v\in H^{1}(\Omega)$, we define
\begin{equation}\label{pb:extopt1}
\Lambda_{m,\delta}(v) = \inf\{
\cQ_{m}(u)\sep u\in V_{v,\delta}\},
\end{equation}
where $\cQ_m(u) = \|\nabla u\|^2_{L^2(\Omega')} + m^2\|u\|^2_{L^2(\Omega')}$
is defined in \eqref{pb:extopt}
and
\begin{equation*}
\begin{split} \mathcal V_{\delta} &= \{\bx\in\Omega'\sep\dist(\bx,\Gamma)<\delta\},\\
V_{v,\delta} &= \bigl\{
u\in H^1(\mathcal V_{\delta},\mathbb{C}^4) \sep u = v\mbox{ on }\Gamma\mbox{ and }u(\bx) =0\mbox{ if }\dist(\bx,\Gamma) = \delta
\bigr\}.
\end{split}
\end{equation*}

\begin{remark}\label{rem.regvoitub}
Note that, since $\Omega$ is a smooth set, there exists $\delta_0>0$ such that, for all $\delta\in(0,\delta_0)$, the set $\cV_\delta$ has the same regularity as $\Omega$.
\end{remark}

\Subsubsection{Existence and uniqueness of minimizers}

\begin{lemma}\label{lem.existence}

For $\delta\in(0,\delta_0)$, $m>0$, and $v\in H^1(\Omega)$,
the minimizers associated with \eqref{pb:extopt} and \eqref{pb:extopt1} exist and are unique.
\end{lemma}
\begin{proof}
Let $(u_n)$ and $(u_{\delta,n})$ be minimizing sequences for $\Lambda_{m}(v)$ and $\Lambda_{m,\delta}(v)$ respectively. These two sequences are uniformly bounded in $H^1$ so that, up to subsequences, they converge weakly to $u\in H^1(\Omega')$ and $v_\delta\in H^1(\cV_\delta)$, respectively. By Rellich-Kondrachov compactness Theorem and the interpolation inequality,
the sequences converges strongly in $H^{s}_\loc$ for any $s\in[0,1)$. The trace theorem ensures then that the convergence also holds in $L^2_\loc(\Gamma)$ and $L^2_\loc(\partial \cV_\delta)$, so that $u\in V_v$ and $u_\delta\in V_{v,\delta}$. Since
\[
\Lambda_m(v) = \lim_{n\to+\infty} \cQ_m(u_n)\geq \cQ_m(u)\geq \Lambda_m(v)
\]
and
\[
\Lambda_{m,\delta}(v) = \lim_{n\to+\infty} \cQ_{m}(u_{\delta,n})\geq \cQ_{m}(u_{\delta,n})\geq \Lambda_{m,\delta}(v),
\]
$u$ and $u_\delta$ are minimizers.

Finally, since $V$ and $V_\delta$ are convex sets and the quadratic form $\cQ_m$ is a strictly convex function, the uniqueness follows.
\end{proof}

\begin{notation}
The unique minimizers associated with $\Lambda_{m}(v)$ and $\Lambda_{m,\delta}(v)$ will be denoted by $u_m(v)$ and $u_{m,\delta}(v)$, respectively, or by $u_m$ and $u_{m,\delta}$ when the dependence on $v$ is clear.
\end{notation}

\subsubsection{Euler-Lagrange equations}

The following lemma gathers some properties related to the Euler-Lagrange equations.

\begin{lemma}\label{lem:ELE}
For all $\delta\in(0,\delta_0)$, $m>0$, and $v\in H^1(\Omega)$, the following holds:
\begin{enumerate}
\item\label{pt:cp1} $(-\Delta+m^2)u_m = 0$ and $(-\Delta+m^2)u_{m,\delta}=0$,
\item\label{pt:cp2} $\Lambda_m(v) = -\braket{\pa_\bn u_m,u_m}_\Gamma$ and $\Lambda_{m,\delta}(v) = -\braket{\pa_\bn u_{m,\delta},u_{m,\delta}}_{\Gamma}$,
\item\label{pt:cp4} $\cQ_m(u) = \Lambda_m(v) + \cQ_m(u-u_m)$ for all $u\in V_v$,
$\cQ_{m}(u) = \Lambda_{m,\delta}(v) + \cQ_{m}(u-u_{m,\delta})$ for all $u\in V_{v,\delta}$,
\end{enumerate}
where $\Lambda_m(v)$ and $V_v$ are defined in \eqref{pb:extopt}, $\Lambda_{m,\delta}(v)$ and $V_{v,\delta}$ are defined in \eqref{pb:extopt1}, and~$\delta_0$ is defined in Remark \ref{rem.regvoitub}.
\end{lemma}

\begin{proof}
Let $v\in H^1_0(\Omega')$. The function
\[
\RR\ni t\mto\cQ_m(u_m + tv)
\]
has a minimum at $t=0$. Hence, the Euler-Lagrange equation is $(-\Delta+m^2)u_m = 0$. The same proof holds for $u_{m,\delta}$. The second point follows from integrations by parts.
And for the last point, let $u\in V_v$. We have, by an integration by parts,
\[\begin{split}
\cQ_m(u-u_m)
&= \cQ_m(u)+\cQ_m(u_m) - 2\Re\braket{u,(-\Delta+m^2)u_m}_{\Omega'} + 2\braket{u_m,\pa_\bn u_m}_\Gamma
\\&
= \cQ_m(u) -\Lambda_m(v),
\end{split}\]
and the result follows. The same proof works for $\Lambda_{m,\delta}(v)$.
\end{proof}

\subsection{Agmon estimates}\label{sec.Agmon}
This section is devoted to the decay properties of the minimizers in the regime $m\to+\infty$.

As an intermediate step, we will need the following localization formulas.

\begin{lemma}\label{lem.IMS}
Let $m>0$ and $\chi$ be any real bounded Lipschitz function on $\Omega'$.
Then,
\begin{equation}\label{eq:idfq1}
\cQ_m(u_m\chi) = -\braket{\pa_\bn u_m,\chi^2u_m}_{\Gamma} + \|(\nabla\chi) u_m \|^2_{L^2(\Omega')}.
\end{equation}
The same holds for $u_{m,\delta}$.
\end{lemma}

\begin{proof}
By definition, we have
\[\begin{split}
\cQ_m(u_m\chi)
&=
m^2\|\chi u_m\|^2_{L^2(\Omega')} + \|(\nabla\chi) u_m + \chi(\nabla u_m)\|^2_{L^2(\Omega')}
\\
&
=
m^2\|\chi u_m\|^2_{L^2(\Omega')}
+ \|(\nabla\chi) u_m \|^2_{L^2(\Omega')} + \| \chi(\nabla u_m)\|^2_{L^2(\Omega')}\\
&\hspace*{5.75cm}+ 2\Re\braket{u_m\chi, \nabla\chi\cdot \nabla u_m}_{\Omega'}.
\end{split}\]
Then, by an integration by parts,
\[\begin{split}
\| \chi(\nabla u_m)\|^2_{L^2(\Omega')} &= -\braket{\pa_\bn u_m,\chi^2u_m}_{\Gamma} -2\Re\braket{u_m\chi, \nabla\chi\cdot \nabla u_m}_{\Omega'}\\
&\hspace*{3.2cm}+ \Re\braket{-\Delta u_m,\chi^2 u_m}_{\Omega'}.
\end{split}\]
It remains to use Lemma \ref{lem:ELE} to get
\[
\cQ_m(u_m\chi) = -\braket{\pa_\bn u_m,\chi^2u_m}_{\Gamma} + \|(\nabla\chi) u_m \|^2_{L^2(\Omega')}.
\]
The conclusion follows.
\end{proof}

We can now establish the following important proposition.
\begin{proposition}\label{lem:agmon1}
Let $\gamma\!\in\!(0,1)$.
There exist $C_1,\,C_2\!>\!0$ such that, for all $\delta\!\in\!(0,\delta_0)$, $m>0$, and $v\in H^1(\Omega)$,
\begin{equation}\label{eq.agmon1a}
\|e^{m\gamma\dist(\cdot, \Gamma)}u_m\|^2_{L^2(\Omega')}\leq C_1 \|u_m\|^2_{L^2(\Omega')},
\end{equation}
and
\begin{equation}\label{eq.agmon1b}
(1 - e^{-\gamma m^{1/2}}C_2m^{-1})\Lambda_{m,m^{-1/2}}(v)\leq\Lambda_{m}(v)\leq\Lambda_{m,\delta}(v).
\end{equation}
Here, $\delta_0$ is defined in Remark \ref{rem.regvoitub}.
\end{proposition}

\begin{proof}
Let us first prove \eqref{eq.agmon1a}. Given $\eps>0$, we define
\begin{align*}
\Phi : \bx &\mto \min(\gamma \dist(\bx, \Gamma),\eps^{-1}),
\\
\chi_m : \bx&\mto e^{m\Phi(\bx)},
\end{align*}
and
\begin{align*}
\xi_1\colon [0,1]&\to [0,1]&\xi_2\colon [0,1]&\to [0,1]\\
r&\mto\frac{1-r}{\sqrt{r^2+(1-r)^2}}
&r&\mto\frac{r}{\sqrt{r^2+(1-r)^2}},
\end{align*}
so that $\xi_1^2 + \xi_2^2 = 1$.
We denote $c = \|\xi_1\|_{L^\infty([0,1])} = \|\xi_2\|_{L^\infty([0,1])}>0$.
Let $R>0$.
Let $\chi_{1,m,R},\, \chi_{2,m,R}$ be the Lipschitz quadratic partition of the unity defined by
\[
\chi_{1,m,R}(\bx) = \begin{cases}
1 &\mbox{ if } \dist(\bx,\Gamma)\leq R/2m,\\
\xi_1(2m/R\ \dist(\bx,\Gamma)-1)&\mbox{ if } R/2m\leq\dist(\bx, \Gamma)\leq R/m,\\
0 &\mbox{ if } \dist(\bx, \Gamma)\geq R/m,
\end{cases}
\]
and
\[
\chi_{2,m,R}(\bx) = \begin{cases}
0 &\mbox{ if } \dist(\bx,\Gamma)\leq R/2m,\\
\xi_2(2m/R\ \dist(\bx,\Gamma)-1)&\mbox{ if } R/2m\leq\dist(\bx, \Gamma)\leq R/m,\\
1 &\mbox{ if } \dist(\bx, \Gamma)\geq R/m.
\end{cases}
\]
We get, for $k\in \{1,2\}$,
\[
\|\nabla \chi_{k,m,R}\|_{L^\infty(\Omega')}\leq \frac{2m c}{R}.
\]

Since $\chi_m$ is a bounded, Lipschitz function and is equal to $1$ on $\Gamma$, we get $u_m\chi_m\in V_v$.
By definition and using \eqref{eq:idfq1}, we get
\[\Lambda_m(v) = \cQ_m(u_m) = -\braket{\pa_\bn u_m,u_m}_{\Gamma} = \cQ_m(u_m\chi_m)
- \|(\nabla\chi_m) u_m \|^2_{L^2(\Omega')}.\]
Then, we use the fact that $\nabla (\chi_{1,m,R}^2+\chi_{2,m,R}^2) = 0$ to get
\begin{multline*}
\cQ_m(u_m) =\cQ_m(u_m\chi_m\chi_{1,m,R})+\cQ_m(u_m\chi_m\chi_{2,m,R})- \|(\nabla\chi_m) u_m \|^2_{L^2(\Omega')}\\
-\|(\nabla\chi_{1,m,R})\chi_m u_m \|^2_{L^2(\Omega')}
-\|(\nabla\chi_{2,m,R})\chi_m u_m \|^2_{L^2(\Omega')}.
\end{multline*}
Since
$\cQ_m(u_m\chi_m\chi_{1,m,R})\geq \Lambda_m(v)$
and
\[
\begin{split}
\cQ_m(u_m\chi_m\chi_{2,m,R})
&\geq m^2\norm{u_m\chi_m\chi_{2,m,R}}_{L^2(\Omega')}^2
\\&
= m^2\norm{u_m\chi_m}_{L^2(\Omega')}^2 - m^2\norm{u_m\chi_m\chi_{1,m,R}}_{L^2(\Omega')}^2,
\end{split}
\]
we get that
\begin{equation*}
\begin{split}
m^2\Big(1-\gamma^2-\frac{8{c}^2}{R^2}\Big)\norm{u_m\chi_m}_{L^2(\Omega')}^2
&\leq
m^2\norm{u_m\chi_m\chi_{1,m,R}}_{L^2(\Omega')}^2
\\
&\leq
m^2 e^{2m\min\left( \sfrac{\gamma R}{m}, \sfrac{1}{\eps}\right)}\norm{u_m}_{L^2(\Omega')}^2\\
&\leq
m^2 e^{2\gamma R}\norm{u_m}_{L^2(\Omega')}^2.
\end{split}
\end{equation*}
Taking $R>0$ big enough so that $1-\gamma^2-\sfrac{8{c}^2}{R^2}>0$, we have
\[
\norm{u_m\chi_m}_{L^2(\Omega')}^2\leq C\norm{u_m}_{L^2(\Omega')}^2,
\]
where $C$ does not depend on $\eps$. Taking the limit $\eps\to0$ and using the Fatou lemma we obtain \eqref{eq.agmon1a}.

Let us now prove \eqref{eq.agmon1b}. We have for every $\delta\in(0,\delta_0)$ that $V_{v,\delta}\subset V_v$, so that
\[
\Lambda_{m}(v)\leq\Lambda_{m,\delta}(v).
\]
Let us consider a Lipschitz function $\tilde \chi_m : \Omega' \to [0,1]$ defined for all $\bx\in \Omega'$ by
\[
\tilde \chi_m(\bx) =
\begin{cases}
1 & \mbox{ if }\dist(\bx,\Gamma)\leq \sfrac{1}{2m^{1/2}},\\
0 & \mbox{ if }\dist(\bx,\Gamma)\geq \sfrac{1}{m^{1/2}},
\end{cases}
\]
with $\|\nabla \tilde\chi_{m}\|_{L^\infty(\Omega')}\leq 2 cm^{1/2}$. Thanks to \eqref{eq:idfq1}, we find
\begin{equation}\label{eq.agmon1a bis}
\Lambda_{m,m^{-1/2}}(v)\leq \cQ_m(u_m\tilde \chi_m) = \Lambda_m(v) + \|u_m\nabla \tilde \chi_m\|^2_{L^2(\Omega')}.
\end{equation}
Then, by \eqref{eq.agmon1a} we have
\[ \|u_m\nabla \tilde \chi_m\|^2_{L^2(\Omega')}\leq e^{-\gamma m^{1/2}}4 c^2m\|e^{m\gamma\dist(\cdot, \Gamma)}u_m\|^2_{L^2(\Omega')}\leq
C_1e^{-\gamma m^{1/2}}4 c^2m \|u_m\|^2_{L^2(\Omega')}.\]
Observing that
\[m\|u_{m}\|^2_{L^2(\Omega')}\leq m^{-1}\Lambda_{m}(v),\]
and using \eqref{eq.agmon1a bis} we easily get \eqref{eq.agmon1b}.
\end{proof}

\subsection{Optimization problem in a tubular neighborhood}\label{sec.tub}
From Proposition \ref{lem:agmon1}, we see that, in order to estimate $\Lambda_{m}(v)$, it is sufficient to estimate $\Lambda_{m,m^{-1/2}}(v)$. For that purpose, we will use tubular coordinates.

\subsubsection{Tubular coordinates}
Let $\iota$ be the canonical embedding of $\Gamma$ in $\RR^3$ and $g$ the induced metric on $\Gamma$. $(\Gamma,g)$ is a $\cC^{2}$ Riemannian manifold, which we orientate according to the ambient space. Let us introduce the map $\Phi:\Gamma\times(0,\delta)\to\cV_{\delta}$ defined by the formula
\begin{equation*}
\Phi(s,t)=\iota(s)+t\bn(s),
\end{equation*}
where $\cV_\delta$ is defined below \eqref{pb:extopt1}.
The transformation $\Phi$ is a $\cC^{1}$ diffeomorphism for all $\delta\in(0,\delta_0)$ provided that $\delta_0$ is sufficiently small. The induced metric on $\Gamma\times(0,\delta)$ is given by
\[G=g\circ (\sfId+tL(s))^{2}+\dx t^2,\]
where $L(s)=d\bn_{s}$ is the second fundamental form of the boundary at $s\in\Gamma$, see Notation \ref{not.def1}.

Let us now describe how our optimization problem is transformed under the change of coordinates. For all $u\in L^{2}(\mathcal V_{\delta})$, we define the pull-back function
\begin{equation}\label{eq:trans-st}
\widetilde u(s,t):= u(\Phi(s,t)).
\end{equation}
For all $u\in H^{1}(\cV_{\delta})$, we have
\begin{equation}\label{eq:bc;n}
\int_{\cV_{\delta}}|u|^{2}\dx\bx=\int_{\Gamma\times(0,\delta)}|\widetilde u(s,t)|^{2}\,\tilde a\dx \Gamma \dx t\,
\end{equation}
and
\begin{equation}\label{eq:bc;qf}
\int_{\cV_{\delta}}|\nabla u|^{2}\dx\bx= \int_{\Gamma\times(0,\delta)} \Big[\langle\nabla_{s} \widetilde u,\tilde g^{-1}\nabla_{s} \widetilde u\rangle +|\partial_{t}\widetilde u|^{2}\Big]\,\tilde a\dx \Gamma\dx t,
\end{equation}
where
\[\tilde g=\big(\sfId+tL(s)\big)^2,\]
and $\tilde a(s,t)= |\tilde g(s,t)|^{\sfrac{1}{2}}$. Here $\langle\cdot\,,\cdot\rangle$ is the Euclidean scalar product and $\nabla_{s}$ is the differential on $\Gamma$ seen through the metric.
Since
$L(s)$ is self-adjoint on $T_s\Gamma$,
we have the exact formula
\begin{equation}\label{eq.Taylor3}
\tilde{a}(s,t) = 1+t\kappa(s) + t^2K(s),
\end{equation}
where $\kappa$ and $K$ are defined in Notation \ref{not.def1}. In the following, we assume that
\begin{equation}\label{eq:depdeltam}
\delta = m^{-1/2}.
\end{equation}
In particular, we will use \eqref{eq:bc;n} and \eqref{eq:bc;qf} with this particular choice of $\delta$.

\Subsubsection{The rescaled transition optimization problem in boundary coordinates}\label{sec.rescaled}

We introduce the rescaling
\[(s,\tau)=(s,mt),\]
and the new weights
\begin{equation}\label{eq:Jac-a'}
\widehat a_m(s,\tau)=\tilde a(s,m^{-1}\tau),\qquad \widehat g_m(s,\tau)=\tilde g(s,m^{-1}\tau).
\end{equation}

\begin{remark}\label{rem.m1}
Note that there exists $m_1\geq1$ such that for all $m\geq m_1$, $s\in \Gamma$ and $\tau\in[0,m^{1/2})$, we have $\widehat a_m(s,\tau)\geq 1/2$.
\end{remark}

We set
\begin{equation}\label{eq:dom-hat}
\begin{split}
\widehat{\mathcal V}_{m}&=\Gamma\times(0,\sqrt{m}),\\
\widehat V_{m}&=\{u\in H^1(\widehat{\mathcal V}_{m}, \CC^4;\widehat a_m \dx \Gamma \dx\tau):~u(\cdot,\sqrt{m})=0\},\\
\widehat{\scrQ}_{m}(u)&=m^{-1}\int_{\widehat{\mathcal V}_{m}}\Bigl(\langle\nabla_{s} u,\widehat g^{-1}_m\nabla_{s} u\rangle+m^2|\partial_\tau u|^2\Bigr)\widehat a_m \dx \Gamma \dx\tau\\[-5pt]
&\hspace*{6cm}+m\int_{\widehat{\mathcal V}_{m}}|u|^2\widehat{a}_m\dx \Gamma\dx \tau,\\
\widehat{\scrL}_m&=-m^{-1} \widehat a^{-1}_m\nabla_s(\widehat a_m \widehat g^{-1}_m\nabla_s)
+m\left(-\widehat a^{-1}_m\partial_\tau\widehat a_m\partial_\tau
+1\right).
\end{split}
\end{equation}
\begin{notation}\label{notnation curvatures}
Given $m\geq m_1$, and
$\kappa,\, K\in\RR$, we define
\begin{align*}
a_{m,\kappa,K} : (0,\sqrt{m})& \to \RR \\[-8pt]
\tau&\mto1 + \frac{\tau\kappa}{m} + \frac{\tau^2K}{m^2}.
\end{align*}
\end{notation}
We let
\begin{equation}\label{eq:bcur}
A =\|\kappa\|_{L^\infty(\Gamma)} \quad\text{and}\quad  B = \|K\|_{L^\infty(\Gamma)}.
\end{equation}
\begin{remark}\label{rem:m0}
We can assume (up to taking a larger $m_1$) that for any
\[
(m,\kappa,K)\in[m_1,+\infty)\times[-A,A]\times [-B,B],
\]
we have
$
a_{m,\kappa,K}(\tau)\geq 1/2
$
for all $\tau\in(0,\sqrt{m})$.
\end{remark}
In the following, we assume that $(m,\kappa,K)\in[m_1,+\infty)\times[-A,A]\times [-B,B]$.

\subsection{One dimensional optimization problem with parameters}
We denote by $\widehat{\scrQ}_{m,\kappa,K}$ the \enquote{transverse} quadratic form defined for $u\in H^1((0,\sqrt{m}), a_{m,\kappa,K}\dx\tau)$ by
\[
\widehat{\scrQ}_{m,\kappa,K}(u)
= \int_{0}^{\sqrt{m}}
\big(|\partial_\tau u|^2 + |u|^2\big)a_{m,\kappa,K} \dx\tau.
\]
We let
\begin{equation}\label{eq:optiextpar}
\Lambda_{m,\kappa,K} = \inf\{
\widehat{\scrQ}_{m,\kappa,K}(u) \sep u\in \widehat{V}_{m,\kappa,K}
\},
\end{equation}
where
\[
\widehat{V}_{m,\kappa,K} = \left\{
u \in H^1((0,\sqrt{m}), a_{m,\kappa,K}\dx\tau)\sep u(0) = 1, \ u(\sqrt{m})=0
\right\}.
\]
The following lemma follows from the same arguments as for Lemma \ref{lem.existence}.
\begin{lemma}\label{unique minim 1d}
There is a unique minimizer $u_{m,\kappa,K}$ for the optimization problem \eqref{eq:optiextpar}.
\end{lemma}

\begin{lemma}\label{lem.ipp1D}
Let $u\in H^2((0,\sqrt{m}), a_{m,\kappa,K}\dx\tau)$ and $v\in H^1((0,\sqrt{m}), a_{m,\kappa,K}\dx\tau)$ {} be such that $u(\sqrt{m}) = v(\sqrt{m})= 0$. We have
\begin{multline}\label{eq:IPPno}
\int_0^{\sqrt{m}}\langle
\pa_\tau u,\pa_\tau v
\rangle\, a_{m,\kappa,K}\dx\tau
+
\int_0^{\sqrt{m}}\langle
u,v
\rangle\, a_{m,\kappa,K}\dx\tau
\\
=
\int_0^{\sqrt{m}}\bigl\langle
\widehat{\cL}_{m,\kappa,K} u,v
\bigr\rangle\, a_{m,\kappa,K}\dx\tau
-\langle\pa_\tau u(0),v(0)\rangle,
\end{multline}
where
\[
\widehat{\cL}_{m,\kappa,K} =- a^{-1}_{m,\kappa,K}\partial_\tau a_{m,\kappa,K}\partial_\tau
+1= -\pa^2_\tau - \frac{m^{-1}\kappa + m^{-2}2K\tau}{1 + m^{-1}\kappa\tau + m^{-2}K\tau^{2}}\pa_\tau + 1.
\]
\end{lemma}
\begin{proof}
The lemma follows essentially by integration by parts and Notation \ref{notnation curvatures}.
\end{proof}
\begin{lemma}\label{lem.EL1D}
We have that $u_{m,\kappa,K}\in \cC^\infty([0,\sqrt{m}])$ and
\[\widehat{\cL}_{m,\kappa,K} u_{m,\kappa,K} = 0,\qquad \Lambda_{m,\kappa,K}=-\pa_\tau u_{m,\kappa,K} (0),\]
where $u_{m,\kappa,K}$ is defined in Lemma \ref{unique minim 1d}.

Moreover, for all $u\in \widehat{V}_{m,\kappa,K}$,
\[\widehat{\scrQ}_{m,\kappa,K}(u) = \Lambda_{m,\kappa,K}+\widehat{\scrQ}_{m,\kappa,K}(u-u_{m,\kappa,K}).\]
\end{lemma}

\begin{proof}
This follows from Lemma \ref{lem.ipp1D}.
\end{proof}

The aim of this section is to establish an accurate estimate of $\Lambda_{m,\kappa,K}$.
\begin{proposition}\label{lem:ext}
There exists a constant $C>0$ such that for all
\[
(m,\kappa,K)\in[m_1,+\infty)\times[-A,A]\times [-B,B],
\]
we have
\[\Bigl|\Lambda_{m,\kappa,K} - \Bigl(1 + \frac{\kappa}{2m} + \frac{1}{m^2}\Bigl(\frac{K}{2}-\frac{\kappa^2}{8}\Bigr)\Bigr)\Bigr|\leq Cm^{-3},\]
and
\[\biggl|\int_0^{\sqrt{m}}|u_{m,\kappa,K}|^2a_{m,\kappa,K}\dx \tau - \frac{1}{2}\biggr| \leq Cm^{-1}.\]
\end{proposition}

\begin{proof}
By Lemmas \ref{unique minim 1d} and \ref{lem.EL1D}, the unique solution $u_{m,\kappa,K}$ of the problem satisfies
\[
\Bigl(
-\pa_\tau^2
-\frac{m^{-1}\kappa + m^{-2}2K\tau}{1 + m^{-1}\kappa \tau + m^{-2}K\tau^2}\pa_\tau
+1
\Bigr)u_{m,\kappa,K} = 0.
\]
We expand formally $u_{m,\kappa,K}$ as $u_0 + m^{-1}u_1 + m^{-2}u_2 + \scrO(m^{-3})$:
\begin{enumerate}
\item For the \emph{zero order} term, we get
\[
(-\pa_\tau^2 +1)u_0 = 0 \quad\text{and}\quad  u_0(1) = 1,\ \lim_{\tau\to \infty} u_0(\tau) = 0,
\]
so that $u_0(\tau)= e^{-\tau}$.
\item At the \emph{first order},
\[
(-\pa_\tau^2 +1)u_1 = \kappa\pa_\tau u_0 = -\kappa e^{-\tau} \quad\text{and}\quad  u_1(1) = 0,\ \lim_{\tau\to \infty} u_1(\tau) = 0,
\]
so that $u_1(\tau)= -\Psfrac{\kappa}{2}\tau e^{-\tau}$.
\item At the \emph{second order},
\[\begin{split}
(-\pa_\tau^2 +1)u_2 &= \kappa\pa_\tau u_1 + (\kappa^2-2K)\tau \pa_\tau u_0 = -\frac{\kappa^2}{2}e^{-\tau} + \Bigl(\frac{3\kappa^2}{2}-2K\Bigr)\tau e^{-\tau},
\\
u_2(0) &= 0 \quad\text{and}\quad \lim_{\tau\to \infty} u_2(\tau) = 0,
\end{split}\]
so that
$u_2(\tau)=
\bigl(
\frac{\kappa^2}{8}-\frac{K}{2}
\bigr)\tau e^{-\tau}
+
\bigl(
\frac{3\kappa^2}{8}-\frac{K}{2}
\bigr)\tau^2e^{-\tau}.
$
\end{enumerate}
This formal construction leads to define a possible approximation of $u_{m,\kappa,K}$.
Consider
\begin{equation}\label{eq:formaltestfun1}\begin{split}
v_{m,\kappa,K}(\tau) :&= \chi_m(\tau)\left(u_0(\tau) + m^{-1}u_1(\tau) + m^{-2}u_2(\tau)\right),\\
\chi_m(\tau) &= \chi(\tau/\sqrt{m}),
\end{split}\end{equation}
where $\chi : \RR_+\mto [0,1]$ is a smooth function such that
\[
\chi(\tau) = \begin{cases}
1 &\mbox{ if }\tau\in[0,1/2],\\
0 &\mbox{ if }\tau\geq 1.
\end{cases}\
\]
In the following, we denote $v_m \equiv v_{m,\kappa,K}$ to shorten the notation.

We immediately get that $v_m$ belongs to $\widehat{V}_{m,\kappa,K}$. Note that
\begin{equation}\label{eq.deriv-vm}
-\partial_{\tau}v_{m}(0)=1 + \frac{\kappa}{2m} + m^{-2}\Bigl(
\frac{K}{2} - \frac{\kappa^2}{8}
\Bigr)
\end{equation}
and
\begin{equation}\label{eq.approx-zero}
\|\widehat{\cL}_{m,\kappa,K} v_m\|_{L^2((0,\sqrt{m}), a_{m,\kappa,K}\dx\tau)} = \scrO(m^{-3}).
\end{equation}
Using Lemmas \ref{lem.ipp1D} and \ref{lem.EL1D}, we have
\[ \Lambda_{m,\kappa,K}=
\int_0^{\sqrt{m}}\big\langle
\pa_\tau u_{m,\kappa,K},\pa_\tau v_m
\big\rangle a_{m,\kappa,K}\dx\tau
+
\int_0^{\sqrt{m}}\big\langle
u_{m,\kappa,K},v_m
\big\rangle a_{m,\kappa,K}\dx\tau\,
\]
and
\[\begin{split}
&
\Lambda_{m,\kappa,K}
=\int_0^{\sqrt{m}}\big\langle
\widehat{\cL}_{m,\kappa,K} v_m,u_{m,\kappa,K}
\big\rangle a_{m,\kappa,K}\dx\tau
-
\pa_\tau v_{m}(0).
\end{split}\]
By Lemma \ref{lem.ipp1D}, the Cauchy-Schwarz inequality, \eqref{eq.deriv-vm}, and \eqref{eq.approx-zero}, we see that
\[\begin{split}
\Bigl|\Lambda_{m,\kappa,K}&-\Bigl(
1 + \frac{\kappa}{2m} + m^{-2}\Bigl(
\frac{K}{2} - \frac{\kappa^2}{8}
\Bigr)
\Bigr)\Bigr|
\\
&= \biggl|
\int_0^{\sqrt{m}}\big\langle
\widehat{\cL}_{m,\kappa,K} v_m,
u_{m,\kappa,K}
\big\rangle a_{m,\kappa,K}\dx\tau
\biggr|
\\
&\leq
\|\widehat{\cL}_{m,\kappa,K} v_m\|_{L^2((0,\sqrt{m}), a_{m,\kappa,K}\dx\tau)}\cdot
\|u_{m,\kappa,K}\|_{L^2((0,\sqrt{m}), a_{m,\kappa,K}\dx\tau)}
\\
&\leq \Lambda_{m,\kappa, K}^{\sfrac{1}{2}} \|\widehat{\cL}_{m,\kappa,K} v_m\|_{L^2((0,\sqrt{m}), a_{m,\kappa,K}\dx\tau)} \\
&\leq Cm^{-3}\Lambda_{m,\kappa, K}^{\sfrac{1}{2}}.
\end{split}\]
From this, it follows first that $\Lambda_{m,\kappa,K}=\scrO(1)$ uniformly in $(\kappa, K)$, and then the first estimate of the proposition is established.
Using Lemmas \ref{lem.ipp1D} and \ref{lem.EL1D}, the fact that $v_{m}(0)-u_{m,\kappa,K}(0) = 0$, and Cauchy-Schwarz inequality, we have
\begin{align*}
\widehat{\scrQ}_{m,\kappa,K}&(v_{m}-u_{m,\kappa,K})\\
&\hspace*{-.5cm}\leq \|\widehat{\cL}_{m,\kappa,K} (v_m-u_{m,\kappa,K})\|_{L^2((0,\sqrt{m}), a_{m,\kappa,K}\dx\tau)}\|v_m-u_{m,\kappa,K}\|_{L^2((0,\sqrt{m}), a_{m,\kappa,K}\dx\tau)}\\
&\hspace*{-.5cm}\leq Cm^{-3}\|v_m-u_{m,\kappa,K}\|_{L^2((0,\sqrt{m}), a_{m,\kappa,K}\dx\tau)}.
\end{align*}
The second estimate of the proposition follows since
\[\|v_m-u_{m,\kappa,K}\|_{L^2((0,\sqrt{m}), a_{m,\kappa,K}\dx\tau)}^2\leq\widehat{\scrQ}_{m,\kappa,K}(v_{m}-u_{m,\kappa,K})\,\]
and $\|v_{m}\|_{L^2((0,\sqrt{m}), a_{m,\kappa,K}\dx\tau)}^2=\Psfrac{1}{2}+\scrO(m^{-1})$.
\end{proof}
\subsection{Asymptotic study of $\Lambda_{m,m^{-1/2}}(v)$.}\label{sec.proof-prop}
From Proposition \ref{lem:ext} and \eqref{eq:dom-hat}, we \hbox{deduce} the following lower bound.
\begin{corollary}\label{cor:estilowreg}

The following holds:
\begin{enumerate}
\item \label{pt.coroextreg2}
Assume that $\Gamma$ is $\cC^2$. There exists $C\!>\!0$ such that, for every $m\!\geq\!m_1$, $v\!\in\!H^1(\Omega)$,
\[
Cm\norm{v}_{H^1(\Omega)}^2
\geq\Lambda_{m,m^{-1/2}}(v)
\geq \left(m\|v\|^2_{L^2(\Gamma)} + \int_\Gamma \frac{\kappa}{2}\,|v|^2\dx \Gamma\right)
-\frac{C}{m}\,\|v\|^2_{L^2(\Gamma)}.
\]
\item \label{pt.coroextreg4} Assume that $\Gamma$ is $\cC^{2,1}$. There exists $C>0$ such that, for every $m\geq m_1$, $v\in H^1(\Omega)$,
\[\begin{split}
\left(m\|v\|^2_{L^2(\Gamma)} + \int_\Gamma \frac{\kappa}{2}\,|v|^2\dx \Gamma\right)+o(1)&\geq\Lambda_{m,m^{-1/2}}(v).
\end{split}\]
Here, the term $o(1)$ depends on $v$ (not only on the $H^1$ norm of $v$).
\end{enumerate}
\end{corollary}
\begin{proof}
By Proposition \ref{lem:ext}, the lower bound of Point \eqref{pt.coroextreg2} follows. Let us focus on Point \eqref{pt.coroextreg4}.

By the extension theorem for Sobolev functions (see for instance \cite[\S 5.4]{evans1998partial}), there exist a constant $C>0$ and, for all $v\in H^1(\Omega)$, a function $Ev\in H^1(\RR^3)$ that extends~$v$ and such that $\norm{Ev}_{H^1(\RR^3)}\leq C\norm{v}_{H^1(\Omega)}$.

Let us define the test function $u_m$ by $u_m = v \widetilde{u}_m$ where\vspace*{-3pt}
\[
\widetilde{u}_m\circ\Phi(s,t)= \begin{cases}
v_{m,\kappa(s),K(s)}(mt)&\mbox{for all } (s,t)\in\Gamma\times[0,{m}^{-1/2}],\\
0&\mbox{for all } (s,t)\in\Gamma\times[{m}^{-1/2},+\infty).
\end{cases}
\]
Here, the function $v_m$ is defined in \eqref{eq:formaltestfun1}.

Let us first prove a general formula. Consider $u\in H^2(\Omega';\RR)$ and $v\in H^1(\Omega';\CC^4)$. With an integration by parts and using the fact that $u$ is real-valued,
\[\begin{split}
\|\nabla(vu)\|^2_{L^2(\Omega')}
&=
\norm{u\nabla v + v\nabla u}_{L^2(\Omega')}^2
\\
&=
\norm{u\nabla v}_{L^2(\Omega')}^2
+\norm{ v\nabla u}_{L^2(\Omega')}^2
+2\Re\braket{u\nabla v,v\nabla u}_{\Omega'}
\\
&=
\norm{u\nabla v}_{L^2(\Omega')}^2
+\Re\braket{uv,-v\Delta u}_{\Omega'}
-\Re\braket{v\pa_\bn u, vu}_{\Gamma}
\\
&=
\norm{u\nabla v}_{L^2(\Omega')}^2
+\braket{uv,-v\Delta u}_{\Omega'}
-\braket{v\pa_\bn u, vu}_{\Gamma}.
\end{split}\]
With an integration by parts only in the tangential direction,
\[\braket{uv,-v\Delta u}_{\Omega'}=\braket{uv,-v\Delta_t u}_{\Omega'}+2\Re\braket{u\nabla_s v,v\nabla_s u}_{\Omega'}+\|u\nabla_s v\|^2_{L^2(\Omega')},\]
where $\nabla_s$ is the tangential derivative and $-\Delta_t$ is the part of the Laplacian involving the second order derivative in the normal variable $t$. Thus, we get\vspace*{-3pt}
\begin{multline*}
\|\nabla(vu)\|^2_{L^2(\Omega')}
=\norm{u\nabla v}_{L^2(\Omega')}^2
+\braket{uv,-v\Delta_t u}_{\Omega'}+2\Re\braket{u\nabla_s v,v\nabla_s u}_{\Omega'}\\
+\|u\nabla_s v\|^2_{L^2(\Omega')}
-\braket{v\pa_\bn u, vu}_{\Gamma}.
\end{multline*}
By density, this formula can be extended to $u$ in $H^2_t$ and $H^1_s$. Therefore, we can replace~$u$ by $\widetilde{u}_m$. We get\vspace*{-3pt}
\begin{multline}\label{eq.sesqui-expression}
\cQ_m(u_m)
=-\braket{v\pa_\bn \widetilde{u}_m, v\widetilde{u}_m}_{\Gamma}+
\norm{\widetilde{u}_m\nabla v}_{L^2(\Omega')}^2
+\braket{\widetilde{u}_mv,v\left(-\Delta_t+m^2\right) \widetilde{u}_m}_{\Omega'}\\
+2\Re\braket{\widetilde{u}_m\nabla_s v,v\nabla_s \widetilde{u}_m}_{\Omega'}+\|\widetilde{u}_m\nabla_s v\|^2_{L^2(\Omega')}.
\end{multline}
With the explicit expression \eqref{eq:formaltestfun1}, we find
\begin{equation}\label{eq.ub1D}
-\braket{v\pa_\bn \widetilde{u}_m, v\widetilde{u}_m}_{\Gamma}
\leq
m\norm{v}_{L^2(\Gamma)}^2 + \int_\Gamma\frac{\kappa}{2}\,|v|^2\dx\Gamma + \frac{C}{m}\,\norm{v}_{L^2(\Gamma)}^2.
\end{equation}
By using the dominated convergence theorem and the explicit expression $\widetilde{u}_m$, we get that the other terms in \eqref{eq.sesqui-expression} go to $0$. Note here that this argument uses \emph{at most} one derivative of the functions $\kappa(\cdot)$ and $K(\cdot)$ (see the definition of $v_{m,k,K}$ in \eqref{eq:formaltestfun1}). That is why we need $\Gamma$ to be $\cC^{2,1}$.

With the definition of $\Lambda_{m,m^{-1/2}} (v)$, we find\vspace*{-3pt}
\[
\Lambda_{m,m^{-1/2}} (v)\leq m\norm{v}_{L^2(\Gamma)}^2+ \int_\Gamma \frac{\kappa}{2}\,|v|^2\dx\Gamma+o(1).
\]

To get the upper bound of Point \eqref{pt.coroextreg2}, we follow the same steps as before except that $v_{m,k,K}$ is replaced by
\[
\tau \mto \chi_m(\tau)u_0(\tau),
\]
in \eqref{eq:formaltestfun1}. In that case, we only need $\Gamma$ to be $\cC^2$.
\end{proof}

Using Proposition \ref{lem:agmon1}, Corollary \ref{cor:estilowreg} proves in particular \eqref{item.i2} and \eqref{item.i4} in Proposition~\ref{prop:ext}. In this section we address the refinement of the lower bound and the corresponding upper bound. From now on, we assume that $\Gamma$ is $\cC^{2,1}$.

\subsubsection{Preliminary lemmas}
Let us state a few elementary lemmas that we will use later.

\begin{lemma}\label{lem:Sob.gamma0}
There exists $C>0$ such that, for all $f,g\in H^{\sfrac{3}{2}}(\Gamma)$, we have
\[\|fg\|_{H^{\sfrac{3}{2}}(\Gamma)}\leq C\|f\|_{H^{\sfrac{3}{2}}(\Gamma)}\|g\|_{H^{\sfrac{3}{2}}(\Gamma)}.\]
\end{lemma}
\begin{proof}
$H^{\sfrac{3}{2}}(\Gamma)$ is a Banach algebra since $\sfrac{3}{2}>\sfrac{\dim\Gamma}{2}=1$.
\end{proof}

\begin{lemma}\label{lem:Sob.gamma1}
There exists $C>0$ such that, for all $f\in H^{\sfrac{3}{2}}(\Gamma)$, we have
\[\|f\|_{H^{\sfrac{1}{2}}(\Gamma)}\leq C\, \|f\|^{\sfrac{1}{2}}_{L^2(\Gamma)}\,\|f\|^{\sfrac{1}{2}}_{H^1(\Gamma)}.\]
\end{lemma}

\begin{lemma}\label{lem:Sob.gamma2}
There exists $C>0$ such that, for all $f\in H^{\sfrac{1}{2}}(\Gamma,T\Gamma)$ and $g\in H^{1}(\Gamma,\CC)$, we have
\[\left|\int_{\Gamma}f\cdot\nabla_{s} g\dx\Gamma\right|\leq C\|f\|_{H^{\sfrac{1}{2}}(\Gamma)}\|g\|_{H^{\sfrac{1}{2}}(\Gamma)}.\]
Here, $T\Gamma$ is the tangent bundle of $\Gamma$.
\end{lemma}

\Subsubsection{Lower and upper bounds}
\begin{notation}\label{not.vkK}
In the following, we define
\begin{align*}
\widehat{\Pi}_m : H^{1}(\Omega; \CC^4) &\to \widehat{V}_m
\\
v &\mto [(s,\tau)\in \widehat{\mathcal V}_{m}\mapsto v(s)u_{m,\kappa(s),K(s)}(\tau)\in \CC^4],
\end{align*}
where $\widehat{V}_m$ and $\widehat{\mathcal V}_{m}$ are given in \eqref{eq:dom-hat}, and $u_{m,\kappa(s),K(s)}$ is defined by Proposition \ref{unique minim 1d} with $\kappa = \kappa(s) $ and $K = K(s)$.
\end{notation}
\begin{lemma}\label{lem.deriv-u}
Assume that $\Gamma$ is $\cC^{2,1}$.
We have, uniformly in $s$,
\[\int_0^{\sqrt{m}}|\nabla_s u_{m,\kappa(\cdot),K(\cdot)}|^2\dx\tau=\scrO(m^{-2}).\]
\end{lemma}
\begin{proof}
Recall from Lemma \ref{lem.EL1D} that
\[\bigl(- a^{-1}_{m,\kappa,K}\partial_\tau a_{m,\kappa,K}\partial_\tau
+1\bigr)u_{m,\kappa, K}=0.\]
Let us take the derivative with respect to $s$:
\[\bigl(- a^{-1}_{m,\kappa,K}\partial_\tau a_{m,\kappa,K}\partial_\tau
+1 \bigr)\nabla_{s} u_{m,\kappa, K}= \bigl[\nabla_{s}\,,\,a^{-1}_{m,\kappa,K}\partial_\tau a_{m,\kappa,K}\partial_\tau \bigr]u_{m,\kappa, K}.\]
Taking the scalar product with $\nabla_{s} u_{m,\kappa, K}$ and integrating by parts by noticing that $\nabla_{s} u_{m,\kappa, K}(0)=0$, we get
\begin{multline*}
\int_{0}^{\sqrt{m}} |\partial_{\tau} \nabla_{s}u_{m,\kappa, K}|^2a_{m,\kappa, K}\dx\tau+\|\nabla_{s} u_{m,\kappa, K}\|^2_{L^2(a_{m,\kappa,K}\dx\tau)}\\
\leq \left|\bigl\langle \bigl[\nabla_{s},a^{-1}_{m,\kappa,K}\partial_\tau a_{m,\kappa,K}\partial_\tau \bigr]u_{m,\kappa, K},\nabla_{s} u_{m,\kappa, K} \bigr\rangle_{L^2(a_{m,\kappa, K}\dx\tau)}\right|.
\end{multline*}
Since
\[
a^{-1}_{m,\kappa,K}\partial_\tau a_{m,\kappa,K}\partial_\tau = \pa_\tau^2 + \Bigl(\frac{\pa_\tau a_{m,\kappa,K}}{a_{m,\kappa,K}}\Bigr)\pa_\tau,
\]
we get
\[\begin{split}
\Bigl[\nabla_{s},a^{-1}_{m,\kappa,K}\partial_\tau a_{m,\kappa,K}\partial_\tau\Bigr]
&= \Bigl[\nabla_{s}, \Bigl(\frac{\pa_\tau a_{m,\kappa,K}}{a_{m,\kappa,K}} \Bigr)\pa_\tau \Bigr]
= \Bigl(\nabla_s\frac{\pa_\tau a_{m,\kappa,K}}{a_{m,\kappa,K}} \Bigr)\pa_\tau
\\&= \Bigl(
\left(\nabla_s \kappa\right)\pa_\kappa\frac{\pa_\tau a_{m,\kappa,K}}{a_{m,\kappa,K}}
+\left(\nabla_s K\right)\pa_K\frac{\pa_\tau a_{m,\kappa,K}}{a_{m,\kappa,K}}
\Bigr)\pa_\tau.
\end{split}\]

By an explicit computation and the Cauchy-Schwarz inequality, we find
\begin{multline*}
\left| \bigl\langle \bigl[\nabla_{s},a^{-1}_{m,\kappa,K}\partial_\tau a_{m,\kappa,K}\partial_\tau \bigr]u_{m,\kappa, K},\nabla_{s} u_{m,\kappa, K} \bigr\rangle_{L^2(a_{m,\kappa, K}\dx\tau)}\right|\\
\leq Cm^{-1}\|\partial_{\tau}u_{m,\kappa,K}\|_{L^2(a_{m,\kappa, K}\dx\tau)}\cdot\|\nabla_{s}u_{m,\kappa, K}\|_{L^2(a_{m,\kappa, K}\dx\tau)}.
\end{multline*}
Since
\[\|\partial_{\tau}u_{m,\kappa,K}\|_{L^2(a_{m,\kappa, K}\dx\tau)}\leq \sqrt{\Lambda_{m,\kappa, K}},\]
we get by Proposition \ref{lem:ext}
\[\int_{0}^{\sqrt{m}} |\partial_{\tau} \nabla_{s}u_{m,\kappa, K}|^2a_{m,\kappa, K}\dx\tau+\|\nabla_{s} u_{m,\kappa, K}\|^2_{L^2(a_{m,\kappa,K}\dx\tau)}\leq Cm^{-2}.\qedhere\]
\end{proof}

Up to taking a larger $m_1$ in Remark \ref{rem.m1}, we get the following result.

\begin{proposition}\label{lem:effop}
Assume that $\Gamma$ is $\cC^{2,1}$.
There exist positive constants $C>0$ and $m_1>0$ such that, for all $m\geq m_1$, and all $v\in H^{2}(\Omega)$, we have
\[
\big|\Lambda_{m,m^{-1/2}}(v)-\widetilde \Lambda_{m}(v)\big|\leq Cm^{-3/2}\|v\|^2_{H^{3/2}(\Gamma)},
\]
where
\[
\widetilde \Lambda_{m}(v) = m\int_\Gamma
|v|^2
\dx\Gamma
+
\int_\Gamma
\frac{\kappa}{2}\,|v|^2
\dx\Gamma
+
m^{-1}\int_\Gamma \Bigl(
\frac{|\nabla_s v|^2}{2}
+
\Bigl(\frac{K}{2}-\frac{\kappa^2}{8} \Bigr)|v|^2
\Bigr)\dx\Gamma.
\]
More precisely, for all $u\in \widehat{V}_m$ such that
\[
u(s,0)= v(s), \mbox{ for all }s\in \Gamma,
\]
we have
\begin{multline*}
\widehat{\scrQ}_m(u)\geq \widetilde \Lambda_{m}(v)
-\frac{C}{m^{3/2}} \|v\|_{H^{3/2}(\Gamma)} ^2
+\frac{m}{2}
\|u-\widehat{\Pi}_m v\|^2_{L^2(\widehat{\cV}_m, \dx\Gamma\dx\tau)}
\\
+\frac{1}{2m}
\|\nabla_s \bigl(u - \widehat{\Pi}_m v\bigr)\|^2_{L^2(\widehat{\cV}_m, \dx\Gamma\dx\tau)},
\end{multline*}
and
\[\begin{split}
\widehat{\scrQ}_m(\widehat{\Pi}_m(v))\leq \widetilde \Lambda_{m}(v)
+Cm^{-3/2} \bigl(\|v\|^2_{L^2(\Gamma)} + \|\nabla_s v\|^2_{L^2(\Gamma)}\bigr).
\end{split}\]
\end{proposition}

\begin{proof}
Let $v\in H^{2}(\Omega)$. First, let us discuss the upper bound. For that purpose, we insert $\widehat{\Pi}_m v$ in the quadratic form:
\begin{equation*}
\widehat{\scrQ}_{m} (\widehat{\Pi}_m v )=
m\int_\Gamma
\widehat{\scrQ}_{m,\kappa(\cdot),K(\cdot)}(\widehat{\Pi}_m v )\dx
\Gamma
+m^{-1}\int_{\widehat{\mathcal V}_{m}}\bigl\langle\nabla_{s} \widehat{\Pi}_m v,\widehat g^{-1}_m\nabla_{s} \widehat{\Pi}_m v \bigr\rangle\,\widehat a_m \dx \Gamma \dx\tau.
\end{equation*}
We have
\[
m\int_\Gamma \widehat{\scrQ}_{m,\kappa(\cdot),K(\cdot)}(\widehat{\Pi}_m v )\dx\Gamma
= m\int_\Gamma |v|^2\Lambda_{m,\kappa(\cdot),K(\cdot)}\dx\Gamma,
\]
and
\[\int_{\widehat{\mathcal V}_{m}} \bigl\langle\nabla_{s} \widehat{\Pi}_m v,\widehat g^{-1}_m\nabla_{s} \widehat{\Pi}_m v \bigr\rangle\,\widehat a_m \dx \Gamma \dx\tau
\leq (1+Cm^{-\sfrac{1}{2}})\int_{\widehat{\mathcal V}_{m}}|\nabla_{s} \widehat{\Pi}_m v|^2 \dx \Gamma \dx\tau.
\]
Moreover, for all $\varepsilon>0$,
\begin{multline*}
\int_{\widehat{\mathcal V}_{m}}|\nabla_{s} \widehat{\Pi}_m v|^2 \dx \Gamma \dx\tau
\leq(1+\varepsilon)\int_\Gamma
|\nabla_s v|^2
\int_0^{\sqrt{m}}
|u_{m,\kappa(\cdot),K(\cdot)}|^2
\dx\tau
\dx\Gamma\\
+(1+\varepsilon^{-1})\int_{\Gamma}|v|^2\int_0^{\sqrt{m}}|\nabla_{s} u_{m,\kappa(\cdot),K(\cdot)}|^2\dx\tau\dx\Gamma.
\end{multline*}
We now recall Lemma \ref{lem.deriv-u}. Choosing $\varepsilon=m^{-1}$ and using Proposition \ref{lem:ext} we get
\[\int_{\widehat{\mathcal V}_{m}}|\nabla_{s} \widehat{\Pi}_m v|^2 \dx \Gamma \dx\tau\leq (1+Cm^{-1})\frac{1}{2}\int_{\Gamma}|\nabla_{s} v|^2\dx\Gamma+Cm^{-1}\|v\|^2_{L^2(\Gamma)}.\]
Therefore,
\[
\widehat{\scrQ}_{m} (\widehat{\Pi}_m v )\leq m\int_\Gamma |v|^2\Lambda_{m,\kappa(\cdot),K(\cdot)}\dx\Gamma+m^{-1}\frac{1+Cm^{-\sfrac{1}{2}}}{2}\int_{\Gamma}|\nabla_{s} v|^2\dx\Gamma+Cm^{-2}\|v\|^2_{L^2(\Gamma)}.
\]
It only remains to use Proposition \ref{lem:ext} to get the desired upper bound.

Let us now discuss the lower bound. Let $u\in \widehat{V}_m$ such that $u = v$ on $\Gamma$.
By~Lem\-ma~\ref{lem.EL1D}, we have\
\[\begin{split}
\widehat{\scrQ}_m(u) &= m\int_\Gamma \widehat{\scrQ}_{m,\kappa(\cdot),K(\cdot)}(u)\dx \Gamma
+m^{-1}\int_{\widehat{\mathcal V}_{m}}\langle\nabla_{s} u,\widehat g^{-1}_m\nabla_{s} u\rangle\widehat a_m \dx \Gamma \dx\tau
\\
&
= m\int_\Gamma |v|^2\Lambda_{m,\kappa(\cdot),K(\cdot)}\dx \Gamma
+m\int_\Gamma \widehat{\scrQ}_{m,\kappa(\cdot),K(\cdot)}(u-\widehat{\Pi}_m v)\dx \Gamma
\\
&
\hspace*{4cm}+m^{-1}\int_{\widehat{\mathcal V}_{m}}\langle\nabla_{s} u,\widehat g^{-1}_m\nabla_{s} u\rangle\widehat a_m \dx \Gamma \dx\tau.
\end{split}\]
Thus,
\begin{multline}\label{eq:lemboundext8}
\widehat{\scrQ}_m(u)\geq m\int_\Gamma |v|^2\Lambda_{m,\kappa(\cdot),K(\cdot)}\dx \Gamma+m(1-Cm^{-\sfrac{1}{2}})\|u-\widehat{\Pi}_mv\|^2_{L^2(\widehat{\cV}_m, \dx\Gamma\dx\tau)} \\
+m^{-1}(1 - Cm^{-\sfrac{1}{2}})\int_{\widehat{\mathcal V}_{m}}|\nabla_s u|^2 \dx \Gamma \dx\tau.
\end{multline}

We have
\[
\nabla_s u
= \widehat{\Pi}_m \nabla_s v + \bigl(\nabla_s u - \widehat{\Pi}_m \nabla_s v\bigr)
=
\widehat{\Pi}_m \nabla_s v + \nabla_s (u-\widehat{\Pi}_m v) +[ \nabla_s,\,\widehat{\Pi}_m ]\,v\,
\]
and
\[
[ \nabla_s,\,\widehat{\Pi}_m ]\,v(s,\tau) = v(s)\nabla_s u_{m,\kappa(s),K(s)}(\tau).
\]
By Lemma \ref{lem.deriv-u}, we obtain
\[
\int_{\widehat{\mathcal V}_{m}}\bigl|[ \nabla_s,\,\widehat{\Pi}_m ]\,v\bigr|^2\dx \Gamma \dx\tau \leq Cm^{-2}\|v\|^2_{L^2(\Gamma)},
\]
and by Young's inequality,
\begin{multline}\label{eq:lemboundext7}
\int_{\widehat{\mathcal V}_{m}}|\nabla_s u|^2 \dx \Gamma \dx\tau
\\\geq
(1-m^{-1})\int_{\widehat{\mathcal V}_{m}}\bigl|\widehat{\Pi}_m \nabla_s v + \nabla_s (u-\widehat{\Pi}_m v)\bigl|^2 \dx \Gamma \dx\tau
-m
\int_{\widehat{\mathcal V}_{m}}\big|[ \nabla_s,\,\widehat{\Pi}_m ]\,v\big|^2 \dx \Gamma \dx\tau
\\\geq
(1-m^{-1})\int_{\widehat{\mathcal V}_{m}}\bigl|\widehat{\Pi}_m \nabla_s v + \nabla_s (u-\widehat{\Pi}_m v)\bigr|^2 \dx \Gamma \dx\tau
-m^{-1}C\|v\|^2_{L^2(\Gamma)}.
\end{multline}
We also have
\begin{multline}\label{eq:lemboundext2}
\int_{\widehat{\mathcal V}_{m}}\bigl|\widehat{\Pi}_m \nabla_s v + \nabla_s (u-\widehat{\Pi}_m v)\bigr|^2 \dx \Gamma \dx\tau
\geq
\int_{\widehat{\mathcal V}_{m}}\bigl|\widehat{\Pi}_m \nabla_s v \bigr|^2 \dx \Gamma \dx\tau
\\
+
\int_{\widehat{\mathcal V}_{m}}\bigl|\nabla_s (u-\widehat{\Pi}_m v)\bigr|^2 \dx \Gamma \dx\tau
-
\biggl|2\Re
\int_{\widehat{\mathcal V}_{m}}\bigl\langle\widehat{\Pi}_m \nabla_s v, \nabla_s (u-\widehat{\Pi}_m v)\bigr\rangle \dx \Gamma \dx\tau\biggr|,
\end{multline}
and
by Lemmas \ref{lem:Sob.gamma2} and \ref{lem:Sob.gamma0},
\begin{equation*}
\biggl|2\Re
\int_{\widehat{\mathcal V}_{m}}\bigl\langle\widehat{\Pi}_m \nabla_s v, \nabla_s (u-\widehat{\Pi}_m v)\bigr\rangle \dx \Gamma \dx\tau\biggr|
\leq C
\|v\|_{H^{3/2}(\Gamma)}
\bigl\|u - \widehat{\Pi}_mv\bigr\|_{H^{1/2}(\widehat{\cV}_m, \dx\Gamma\dx\tau)}.
\end{equation*}
Then, using Lemma \ref{lem:Sob.gamma1}, we get, for all $\varepsilon_0>0$,
\begin{equation}\label{eq:lemboundext3}
\begin{split}
\biggl|2\Re
\int_{\widehat{\mathcal V}_{m}}\bigl\langle&\widehat{\Pi}_m \nabla_s v, \nabla_s (u-\widehat{\Pi}_m v)\bigr\rangle \dx \Gamma \dx\tau\biggr|\\
&\leq Cm^{-1}\varepsilon_0^{-1}\|v\|_{H^{3/2}(\Gamma)} ^2
+ m^{2}\varepsilon_0\bigl\|u - \widehat{\Pi}_mv\bigr\|^2_{L^2(\widehat{\cV}_m, \dx\Gamma\dx\tau)}\\
&\hspace*{3.55cm}+\varepsilon_0\bigl\|u - \widehat{\Pi}_mv\bigr\|^2_{H^1(\widehat{\cV}_m, \dx\Gamma\dx\tau)}
\\
&\leq Cm^{-1}\varepsilon_0^{-1}\|v\|_{H^{3/2}(\Gamma)} ^2
+ (m^{2}+1)\varepsilon_0\bigl\|u - \widehat{\Pi}_mv\bigr\|^2_{L^2(\widehat{\cV}_m, \dx\Gamma\dx\tau)}\\
&\hspace*{3.55cm}+\varepsilon_0\bigl\|\nabla_s\bigl(u - \widehat{\Pi}_mv\bigr)\bigr\|^2_{L^2(\widehat{\cV}_m, \dx\Gamma\dx\tau)}.
\end{split}\end{equation}
Combining Proposition \ref{lem:ext}, \eqref{eq:lemboundext8}, \eqref{eq:lemboundext7}, \eqref{eq:lemboundext2} and \eqref{eq:lemboundext3}, we finally obtain
\[\begin{split}
\widehat{\scrQ}_m(u)\geq
m\int_\Gamma
|v|^2
&\dx\Gamma
+
\int_\Gamma
\frac{\kappa}{2}\,|v|^2
\dx\Gamma
+
m^{-1}\int_\Gamma \Bigl(
\frac{|\nabla_s v|^2}{2}
+
\Bigl(\frac{K}{2}-\frac{\kappa^2}{8}\Bigr)|v|^2
\Bigr)\dx\Gamma
\\&
-C(m^{-2}+\varepsilon_0^{-1}m^{-2}+m^{-3/2}) \|v\|_{H^{3/2}(\Gamma)} ^2
\\&
+m(1-Cm^{-1/2})\left(1-\varepsilon_0-\varepsilon_0 m^{-2}\right)
\bigl\|u-\widehat{\Pi}_mv\bigr\|^2_{L^2(\widehat{\cV}_m, \dx\Gamma\dx\tau)}
\\
&
+m^{-1}(1-Cm^{-1/2})\left(1-\varepsilon_0\right)
\bigl\|\nabla_s \bigl(u - \widehat{\Pi}_mv\bigr)\bigr\|^2_{L^2(\widehat{\cV}_m, \dx\Gamma\dx\tau)}.
\end{split}\]
Taking $\varepsilon_0 = 3/4$ and $m$ large enough, we get the result.
\end{proof}

\subsection{End of the proof of Proposition \ref{prop:ext}}
Item \eqref{item.ii} of Proposition \ref{prop:ext} follows from Propositions \ref{lem:effop} and \ref{lem:agmon1}. It only remains to prove \eqref{item.iii}. Consider the minimizer $u_{m}$ and a cut off function $\chi_{m}$ supported in a neighborhood of width $m^{-\sfrac{1}{2}}$ near the boundary. Then, we set
\[\check u_{m}(s,\tau)= (\chi_{m} u_{m})\circ \Phi(s,m^{-1}\tau).\]
We now use the lower bound in Proposition \ref{lem:effop}, that is,
\[
\widehat{\scrQ}_m(\check u_{m})\geq \widetilde \Lambda_{m}(v)
+\frac{m}{2}\,
\bigl\|\check u_{m}-\widehat{\Pi}_m v\bigr\|^2_{L^2(\widehat{\cV}_m, \dx\Gamma\dx\tau)}-\frac{C}{m^{3/2}}\, \|v\|_{H^{3/2}(\Gamma)} ^2.
\]
Arguing as in the proof of Lemma \ref{lem.IMS} and recalling Item \eqref{pt:cp2} in Lemma \ref{lem:ELE}, we get
\[\widehat{\scrQ}_m(\check u_{m})=\cQ_{m, m^{-\sfrac{1}{2}}}(\chi_{m} u_{m})=\widetilde \Lambda_{m}(v)+\|(\nabla \chi_{m}) u_{m}\|^2=(1+\scrO(e^{-cm^{\sfrac{1}{2}}}))\widetilde \Lambda_{m}(v),\]
where we also used \eqref{eq.agmon1a}.
Therefore
\[\bigl\|\check u_{m}-\widehat{\Pi}_m v\bigr\|^2_{L^2(\widehat{\cV}_m, \dx\Gamma\dx\tau)}\leq \frac{C}{m^{5/2}}\, \|v\|_{H^{3/2}(\Gamma)} ^2,\]
and then
\[\left|\|\check u_{m}\|_{L^2(\widehat{\cV}_m, \dx\Gamma\dx\tau)}-\bigl\|\widehat{\Pi}_m v\bigr\|_{L^2(\widehat{\cV}_m, \dx\Gamma\dx\tau)}\right|\leq \frac{C}{m^{5/4}}\, \|v\|_{H^{3/2}(\Gamma)}.\]
Using Proposition \ref{lem:ext}, we get that
\[\Bigl|\bigl\|\widehat{\Pi}_m v\bigr\|^2_{L^2(\widehat{\cV}_m, \dx\Gamma\dx\tau)}-\frac{\|v\|^2_{L^2(\Gamma)}}{2}\Bigr|\leq Cm^{-1}\|v\|^2_{L^2(\Gamma)}.\]
Therefore
\[\Bigl |m\|\chi_{m} u_{m}\|^2_{L^2(\cV_m, \dx\bx)}-\frac{\|v\|^2_{L^2(\Gamma)}}{2} \Bigr|\leq Cm^{-1} \|v\|^2_{H^{3/2}(\Gamma)}.\]
Finally, Item \eqref{item.iii} follows by removing $\chi_{m}$ thanks to \eqref{eq.agmon1a}. The proof of Proposition \ref{prop:ext} is complete.\qed

\section{A vectorial Laplacian with Robin-type boundary conditions}\label{sec.3}
In this section, we study the vectorial Laplacian $L_m^\intt$ associated with the quadratic form $\cQ_m^\intt$ defined in Section \ref{sec.robintro}.

\subsection{Preliminaries: proof of Lemma \ref{lem:sarobint}}\label{sec.Lintdef}
We recall that the domain of $L_m^\intt$ is the set of functions $u\in H^1(\Omega; \CC^4)$ such that the linear application
\[ H^1(\Omega; \CC^4)\ni v\mto \cQ_m^\intt(v,u)\in\CC\]
is continuous for the $L^2$-norm. Using the Green-Riemann formula, we get that the domain is indeed given by
\[\{u\in H^1(\Omega; \CC^4) \sep -\Delta u\in L^2(\Omega; \CC^4),\; (\pa_\bn + \kappa/2+m_0 + 2m\Xi^{-})u=0 \mbox{ on } \Gamma\}.\]
By a classical regularity theorem, we deduce that the domain is included in $H^2(\Omega;\CC^4)$.
The compactness of the resolvent and the discreteness of the spectrum immediately follow.
\subsection{Asymptotics of the eigenvalues}
In this section, we describe the first terms in the asymptotic expansion of the eigenvalues of $L_m^\intt$. This is the aim of the following proposition.
\begin{proposition}\label{lem:asym_eigen_int}
The following properties hold:
\begin{enumerate}
\newcounter{saveenum}
\item \label{pt:conv_int_asym1} For every $k\in \NN^*$, we have $\lim_{m\to +\infty} \lambda_{k,m}^\intt = \lambda_k^2$ where the $(\lambda_k)_{k\in \NN^*}$ are the singular values of $|H^\Omega|$.
\setcounter{saveenum}{\value{enumi}}
\end{enumerate}
Let $\lambda$ be an eigenvalue of $|H^\Omega|$ of multiplicity $k_1\in \NN^*$. Let $k_0\in \NN$ be such that $\lambda_{k_0+k} = \lambda$ for all $k\in\{1,\dots,k_1\}$.
\begin{enumerate}
\setcounter{enumi}{\value{saveenum}}
\item\label{pt:conv_int_asym2} For all $k\in\{1,2,\dots, k_1\}$, we have\vspace*{-3pt}
\[
\lambda_{k_0+k,m}^\intt = \lambda^2 +\frac{\mu_{\lambda,k}}{m} + o\left(\sfrac{1}{m}\right),
\]
where\vspace*{-4pt}
\begin{equation}
\mu_{\lambda,k} :=
\inf_{\substack{V\subset \ker(|H^\Omega|-\lambda)\\
\dim V = k}}
\sup_{\substack{v\in V\\
\norm{v}_{L^2(\Omega)}=1}}
-\frac{\norm{(\pa_\bn+\kappa/2+m_0)v}_{L^2(\Gamma)}^2}{2}.
\end{equation}
\item \label{pt:conv_int_asym3}Let $(u_{k_0+1},\dots,u_{k_0+k_1})$ be an $H^1$-weak limit, when $m\to+\infty$, of a sequence \[(u_{k_0+1,m},\dots,u_{k_0+k_1,m})_{m>0}\] of $L^2$-orthonormal eigenvectors of $L^\intt_m$ associated with the eigenvalues \[(\lambda_{k_0+1,m}^\intt,\dots,\lambda_{k_0+k_1,m}^\intt).\]
Then, for all $v\in\ker(|H^\Omega|-\lambda)$, we have\vspace*{-4pt}
\[-\frac{1}{2}\norm{(\pa_\bn+\kappa/2+m_0)v}_{L^2(\Gamma)}^2=\sum_{k=1}^{k_1}|{\braket{ v, u_{k_0+k}}_\Omega}|^2\mu_{\lambda,k}.\]
\end{enumerate}
Here, $(\lambda_k)_{k\in \NN^*}$ is defined in Notation \ref{not.lambda} and $(\lambda_{k,m}^\intt)_{k\in \NN^*}$ in Notation \ref{not.lambdaint}.
\end{proposition}
For the sake of clarity, we will divide the proof of this proposition in different parts. This will be done in the next section.
\subsection{Proof of Proposition \ref{lem:asym_eigen_int}}

Since $\Dom(H^\Omega)\subset \Dom(\cQ_{m}^\intt)$, we have
\begin{equation}\label{eq:upper_bound_op_int}
\lambda_k^2\geq \lambda_{k,m}^\intt
\end{equation}
for all $k\in \NN^*$ and all $m>0$.

\Subsubsection{Lower bounds}
\begin{lemma}\label{lem:induc1}
Let $k\in\NN$.
The following properties hold:
\begin{enumerate}
\item\label{pt:ind_int_asymp1} For all $j\in \{1,2,\dots, k\}$, we have $\lim_{m\to +\infty} \lambda_{j,m}^\intt = \lambda_j^2$.
\item\label{pt:ind_int_asymp2} For all subsequence $(m_n)_{n\in \NN^*}$ going to $+\infty$ as $n\!\to\!+\infty$, and all $L^2$\nobreakdash-orthonormal family of eigenvectors $(u_{1,m_n}, \dots, u_{k,m_n})$ of $L_{m_n}^\intt$ associated with $(\lambda_{1,m_n}^\intt,\dots, \lambda_{k,m_n}^\intt)$ such that the sequence $(u_{1,m_n}, \dots, u_{k,m_n})_{n\in \NN^*}$ converges weakly in $H^1$, we have that the sequence $(u_{1,m_n}, \dots, u_{k,m_n})_{n\in \NN^*}$ converges strongly in $H^1$ and\vspace*{-3pt}
\begin{equation}\label{eq.limbnd}
\lim_{n\to+\infty}m_{n}\norm{\Xi^-u_{j,m_n}}_{L^2(\Gamma)}^2 = 0
\end{equation}
for all $j\in\{1,\dots,k\}$.
\end{enumerate}
\end{lemma}
\begin{proof}
Let us prove \eqref{pt:ind_int_asymp1} and \eqref{pt:ind_int_asymp2} by induction on $k\in \NN^*$.
\subsubsection*{Case $k\!=\!0$} There is nothing to prove.

\subsubsection*{Case $k\!>\!0$} Assume that \eqref{pt:ind_int_asymp1} and \eqref{pt:ind_int_asymp2} are valid for some $k\!\in\!\NN$. Let $(u_{1,m}, \dots, u_{k+1,m})$ be an $L^2$-orthonormal family of eigenvectors of $L_{m}^\intt$ associated with\hspace{-1.2pt} $(\lambda_{1,m}^\intt,\dots, \lambda_{k+1,m}^\intt)$. By \eqref{eq:upper_bound_op_int} and the trace Theorem \cite[\S 5.5]{evans1998partial}, the sequence $(u_{1,m}, \dots, u_{k+1,m})_{m>0}$ is bounded in $H^1(\Omega; \CC^4)^{k+1}$, and
\begin{equation}\label{eq:ineq_asymp_eigen_int}
\lambda^2_{k+1}\geq \limsup_{m\to+\infty}\lambda_{k+1,m}^\intt\geq \liminf_{m\to+\infty}\lambda_{k+1,m}^\intt.
\end{equation}
Hence there exists a subsequence $(m_n)_{n\in \NN^*}$ going to $+\infty$ as $n\to+\infty$ such that
\[
\lim_{n\to+\infty}\lambda_{k+1,m_n}^\intt = \liminf_{m\to+\infty}\lambda_{k+1,m}^\intt
\]
and $(u_{1,m_n}, \dots, u_{k+1,m_n})_{n\in \NN^*}$ converges weakly in $H^1(\Omega;\CC^4)$ to $(u_1,\dots, u_{k+1})$.

Using the induction assumption, we get that $(u_{1,m_n}, \dots, u_{k,m_n})_{n\in \NN^*}$ converges strongly in $H^1(\Omega;\CC^4)$ to $(u_1,\dots, u_{k})$, $\lim_{m\to +\infty} \lambda_{j,m}^\intt = \lambda_j^2$ and
\[
\lim_{n\to+\infty}m\norm{\Xi^-u_{j,m_n}}_{L^2(\Gamma)}^2 = 0
\]
for all $j\in\{1,\dots,k\}$.
By Rellich-Kondrachov Theorem \cite[\S 5.7]{evans1998partial}, the sequence $(u_{k+1,m_n})$ converges strongly in $L^2(\Omega;\CC^4)$. This shows that $(u_1, \dots,u_{k+1})$ is an $L^2$-orthonormal family. In addition, for all $j_1,j_2\in \{1,\dots, k+1\}$, $j_1\ne j_2$, and all $n\in \NN^*$, we have
\begin{multline*}
0 = \Re\braket{\nabla u_{j_1,m_n},\nabla u_{j_2,m_n}}_\Omega+m_0^2\Re\braket{u_{j_1,m_n},u_{j_2,m_n}}_\Omega
\\
+ \Re\braket{(\kappa/2+m_0) u_{j_1,m_n}, u_{j_2,m_n}}_\Gamma + 2m_{n}\Re\braket{\Xi^- u_{j_1,m_n}, \Xi^-u_{j_2,m_n}}_\Gamma,
\end{multline*}
and taking the limit $n\to+\infty$,
\[
0 = \Re\braket{\nabla u_{j_1},\nabla u_{j_2}}_\Omega +m_0^2\Re\braket{u_{j_1},u_{j_2}}_\Omega+ \Re\braket{(\kappa/2 +m_0)u_{j_1}, u_{j_2}}_\Gamma.
\]
Since
\[
\lim_{n\to+\infty} \cQ_{m_n}^\intt(u_{j,m_n}) = \lambda^2_j = \cQ^\intt(u_{j})
\]
for all $j\in\{1,\dots,k\}$, where $\cQ^\intt$ is defined in \eqref{eq:quadMIT}, we deduce that the $(u_{j})_{1\leq j\leq k}$ are normalized eigenfunctions associated with $(\lambda^2_{j})_{1\leq j\leq k}$. By the min-max theorem, we deduce that
\[
\liminf_{n\to+\infty} \cQ_{m_n}^\intt(u_{k+1,m_n})\geq \cQ^\intt(u_{k+1})
\geq \lambda^2_{k+1}.
\]
Therefore
\[\lim_{m\to+\infty}\lambda_{k+1,m}^\intt = \lambda^2_{k+1}\,\]
and \[
\lim_{n\to+\infty}\norm{\nabla u_{k+1,m_n}}_{L^2(\Omega)} = \norm{\nabla u_{k+1}}_{L^2(\Omega)},
\]
and the strong convergence follows. Note that $\lim_{m\to+\infty}\lambda_{k+1,m}^\intt = \lambda^2_{k+1}$ implies that the previous arguments are valid for every weakly converging subsequence, thus Items~\eqref{pt:ind_int_asymp1} and \eqref{pt:ind_int_asymp2} follow for $k+1$.
\end{proof}

\subsubsection{A technical lemma}
The following lemma is essential in the proof of Items~\eqref{pt:conv_int_asym2} and \eqref{pt:conv_int_asym3}.\enlargethispage{-2\baselineskip}%
\begin{lemma}\label{lem:asympt_intlem1}
Let $k\in \NN^*$ and $m>0$.
Let $u$ \resp  $u_{k,m}$ be a $L^2$-normalized eigenfunction of $|H^\Omega|$ \resp  $L_m^\intt$ associated with the eigenvalues $\lambda$ \resp  $\lambda_{k,m}^\intt$. Then
\begin{equation}\label{eq:iden_asym_int}
m(\lambda_{k,m}^\intt-\lambda^2)\braket{u_{k,m},u}_{\Omega} = -1/2\braket{(\pa_\bn + \kappa/2+m_0)u_{k,m},(\pa_\bn + \kappa/2+m_0) u}_\Gamma.
\end{equation}
\end{lemma}

\begin{proof}
On one hand, that $u\in \Dom(|H^\Omega|)$ yields
$\Xi^- u = 0$ on $\Gamma$. Moreover, since $u\in H^1(\Omega; \CC^4)$ is an eigenfunction of $|H^\Omega|$, we indeed have $u\in \Dom((H^\Omega)^2)$, which means that the linear application
$$ H^1(\Omega; \CC^4)\ni v\mto \braket{H^\Omega u,H^\Omega v}_\Omega\in\CC$$is continuous for the $L^2$-norm. Using the Green-Riemann formula, we then get
$$\Xi^+ (\pa_\bn + \kappa/2+m_0)u = 0\,\mbox{ on }\Gamma.$$ On the other hand, from \eqref{eq:laplacian_int} we have
\[
\begin{cases}
\Xi^+ (\pa_\bn + \kappa/2+m_0)u_{k,m} = 0\\
\Xi^-(\pa_\bn + \kappa/2+m_0+2m)u_{k,m} = 0
\end{cases}\quad\mbox{on }\Gamma.
\]
By an integration by parts, we get
\begin{align*}
(\lambda_{k,m}^\intt-\lambda^2)\braket{u_{k,m},u}_{\Omega}
&= \braket{(-\Delta+m_0^2) u_{k,m},u}_\Omega - \braket{u_{k,m},(-\Delta+m_0^2) u}_\Omega
\\&= -\braket{\pa_\bn u_{k,m},u}_\Gamma +\braket{u_{k,m},\pa_\bn u}_\Gamma
\\&= -\braket{(\pa_\bn+\kappa/2+m_0) u_{k,m},u}_\Gamma +\braket{u_{k,m},(\pa_\bn + \kappa/2+m_0) u}_\Gamma
\\&= \braket{\Xi^-u_{k,m},\Xi^-(\pa_\bn + \kappa/2+m_0) u}_\Gamma
\\&= -1/2m\braket{\Xi^-(\pa_\bn + \kappa/2+m_0)u_{k,m},\Xi^-(\pa_\bn + \kappa/2+m_0) u}_\Gamma.
\end{align*}
\end{proof}

\subsubsection{Proof of Items \eqref{pt:conv_int_asym2}\, and \eqref{pt:conv_int_asym3}}
Let $(u_{1,m_n},\dots, u_{k_0+k_1,m_n})_{n\in \NN^*}$ be a sequence of $L^2$-orthonormal eigenvectors of $L_{m_n}^\intt$ that converges strongly in $H^1(\Omega; \CC^4)^{k_0+k_1}$ to an $L^2$-orthonormal family $(u_1,\dots, u_{k_0+k_1})$ of eigenvectors of $|H^\Omega|$.
We have
\[
\Span(u_{k_0+1}, \dots, u_{k_0+k_1}) = \ker(|H^\Omega|-\lambda).
\]
By \eqref{eq:iden_asym_int},
for all $v=\sum_{k=1}^{k_{1}} a_{k}u_{k_{0}+k}$ we have
\[\begin{split}
-1/2&\norm{(\pa_\bn+\kappa/2+m_0)v}_{L^2(\Gamma)}^2\\[-3pt]
& = -1/2\sum_{k,j=1}^{k_1}\overline{a_k}a_j\braket{(\pa_\bn+\kappa/2+m_0)u_{k_0+k}, (\pa_\bn+\kappa/2+m_0)u_{k_0+j}}_\Gamma
\\[-5pt]&
=\lim_{n\to+\infty}-1/2\sum_{k,j=1}^{k_1}\overline{a_k}a_j\braket{(\pa_\bn+\kappa/2+m_0)u_{k_0+k,m_n}, (\pa_\bn+\kappa/2+m_0)u_{k_0+j}}_\Gamma
\\[-3pt]&
=\lim_{n\to+\infty}\sum_{k,j=1}^{k_1}\overline{a_k}a_jm_n(\lambda^\intt_{k_0+k,m_n}-\lambda^2)\braket{u_{k_0+k,m_n},u_{k_0+j}}_\Omega
\\[-3pt]&
=\lim_{n\to+\infty}\sum_{k,j=1}^{k_1}\overline{a_k}a_jm_n(\lambda^\intt_{k_0+k,m_n}-\lambda^2)\braket{u_{k_0+k},u_{k_0+j}}_\Omega
\\[-3pt]&
=\lim_{n\to+\infty}\sum_{k=1}^{k_1}|a_k|^2m_n(\lambda^\intt_{k_0+k,m_n}-\lambda^2).
\end{split}\]

Since $\lambda^\intt_{k_0+1,m_n}\leq\dots\leq \lambda^\intt_{k_0+k_1,m_n}$ and $\CC^{k_1}$ is finite dimensional,
we get for $j\in\{1,\dots,k_1\}$,
\begin{equation*}
\begin{split}
\lim_{n\to+\infty}m_n(\lambda^\intt_{k_0+j,m_n}-\lambda^2)
&
=
\lim_{n\to+\infty}
\biggl(
\min_{\substack{V\subset \CC^{k_1}\\
\dim V = j}}\;
\max_{\substack{a\in V\\
\norm{a}_{l^2}=1}}
\sum_{k=1}^{k_1}|a_k|^2m_n(\lambda^\intt_{k_0+k,m_n}-\lambda^2)\biggr)
\\[-3pt]&
=
\min_{\substack{V\subset \CC^{k_1}\\
\dim V = j}}\;
\max_{\substack{a\in V\\
\norm{a}_{l^2}=1}}
\lim_{n\to+\infty}\sum_{k=1}^{k_1}|a_k|^2m_n(\lambda^\intt_{k_0+k,m_n}-\lambda^2)
\\&
=
-1/2\min_{\substack{V\subset \CC^{k_1}\\
\dim V = j}}\;
\max_{\substack{a\in V\\
\norm{a}_{l^2}=1}}
\biggl\|(\pa_\bn+\kappa/2+m_0)\sum_{k=1}^{k_{1}} a_{k}u_{k_{0}+k}\biggr\|_{L^2(\Gamma)}^2
\\&=
\inf_{\substack{V\subset \ker(|H^\Omega|-\lambda)\\
\dim V = j}}\;
\sup_{\substack{v\in V\\
\norm{v}_{L^2(\Omega)}=1}}
-\frac{\norm{(\pa_\bn+\kappa/2+m_0)v}_{L^2(\Gamma)}^2}{2} \\&= \mu_{\lambda,j},
\end{split}
\end{equation*}
where $\|(a_1,a_2,\dots,a_{k_1})\|_{l^2}^2 = \sum_{k=1}^{k_1}|a_k|^2$ for all $(a_1,a_2,\dots,a_{k_1})\in \CC^{k_1}$.

We obtain
\[
\lim_{m\to+\infty}m(\lambda^\intt_{k_0+j,m}-\lambda^2) = \mu_{\lambda,j}.
\]

Note that a permutation of the limit and the summation sign at the third line of the calculation above ensures that $(u_{k_0+1},\dots,u_{k_0+k_1})$ is an orthogonal family for the quadratic form
\[
v\mto -\frac{\norm{(\pa_\bn+\kappa/2+m_0)v}_{L^2(\Gamma)}^2}{2}.
\]

This finishes the proof of Proposition \ref{lem:asym_eigen_int}.\qed

\section{Proof of the main theorem}\label{sec.4}
We are now ready to address the proof of Theorem \ref{thm:main}. For the sake of readability, we will divide it in several parts.
\subsection{First term in the asymptotic}
In this part, we work in the energy space without using any regularity result such as Lemma \ref{lem:regu}.
\subsubsection{Upper bound}
Let $K\in\NN^*$ and $( \varphi_{1},\dots, \varphi_{K})$ be an $L^2$-orthonormal family of eigenvectors of $|H^\Omega|$ associated with the eigenvalues $(\lambda_{1},\dots,\lambda_{K})$. Using Proposition~\ref{prop:ext}, we extend these functions outside $\Omega$ by
\[
\widetilde u_{j,m} = \begin{cases}
\varphi_{j}& \mbox{ on } \Omega,\\
u_{m+m_0}(\varphi_{j}) &\mbox{ on } \Omega',
\end{cases}
\]
for $j\in\{1,\dots,K\}$. By Proposition \ref{prop:ext}, we get that
\[
\norm{\widetilde u_{j,m}}_{L^2(\Omega')}^2\leq (m+m_0)^{-2}\Lambda_{m+m_0}(\varphi_{j})\leq \frac{C}{m+m_0},
\]
so that $\widetilde u_{1,m},\dots,\widetilde u_{K,m}$ are linearly independent vectors.
Let $a_1,\dots,a_K\in\CC$, and we denote $\varphi_{m}^a := \sum_{j=1}^Ka_j\widetilde{u}_{j,m}$.
By Lemma \ref{lem:quadf1} and Proposition \ref{prop:ext}, we have
\[\begin{split}
\norm{H_m\varphi_m^a}_{L^2(\RR^3)}^2
&= \norm{\nabla\varphi_m^a}_{L^2(\Omega)}^2 + m_0^2\norm{\varphi_m^a}_{L^2(\Omega)}^2
-m\Re\langle\cB\varphi_m^a, \varphi_m^a\rangle_{\Gamma} + \Lambda_{m+m_0}(\varphi_m^a)
\\
&\leq
\cQ^\intt\biggl(\sum_{j=1}^Ka_j\varphi_{j}\biggr)+o(1) = \sum_{j=1}^K|a_j|^2\lambda_j^2+o(1)\leq \lambda_K^2 \sum_{j=1}^K|a_j|^2+o(1).
\end{split}\]
We deduce that
\begin{equation}\label{eq:bsfirsttot}
\limsup_{m\to+\infty}\lambda_{K,m}^2
\leq \limsup_{m\to+\infty}\;
\sup_{\substack{\varphi_m^a\in\Span(\widetilde{u}_{1,m},\dots,\widetilde{u}_{K,m})
\\
\norm{\varphi_m^a}_{L^2(\RR^3)}=1}}\norm{H_m\varphi_m^a}_{L^2(\RR^3)}^2\leq \lambda_K^2.
\end{equation}
\subsubsection{Lower bound and convergence}
For $m\!\geq\!m_1$, let $K\!\in\!\NN^*$ and $( \varphi_{1,m},\dots, \varphi_{K,m})$ be an $L^2$-orthonormal family of eigenvectors of $|H_m|$ associated with the eigenvalues
$(\lambda_{1,m},\dots,\lambda_{K,m})$. Here, $m_1$ is defined in Remark \ref{rem.m1} and Proposition \ref{lem:effop}.
By \eqref{eq:bsfirsttot}, there exists $C>0$ such that
\begin{equation}\label{eq:unifbound}
C\geq \sup_{\substack{k\in\{1,\dots,K\}\\ m\geq m_1}}\norm{H_m \varphi_{k,m}}_{L^2(\RR^3)}^2.
\end{equation}
Using \eqref{eq:quadIM} and Proposition \ref{prop:ext}, we get, for all $k\in\{1,\dots,K\}$ and all $m\geq m_1$, that
\begin{equation}\label{eq:firstest}\begin{split}
\lambda_{k,m}^2 &= \|H_{m} \varphi_{k,m}\|_{L^2(\RR^3)}^2
\\
&
=
\norm{\nabla \varphi_{k,m}}_{L^2(\Omega)}^2 + m_0^2\norm{ \varphi_{k,m}}_{L^2(\Omega)}^2
-m\braket{\cB \varphi_{k,m}, \varphi_{k,m}}_\Gamma
\\
&\hspace*{2.55cm}+
\Lambda_{m+m_0}( \varphi_{k,m})
+\cQ_{m+m_0}( \varphi_{k,m} - u_{m+m_0}( \varphi_{k,m}))
\\
&
\geq \cQ^\intt_m( \varphi_{k,m})
+(m+m_0)^2\norm{ \varphi_{k,m} - u_{m+m_0}( \varphi_{k,m})}_{L^2(\Omega')}^2\\
&\hspace*{5.05cm}-\frac{C}{m}\norm{ \varphi_{k,m}}_{L^2(\Gamma)}^2.
\end{split}\end{equation}
By the trace theorem, we deduce that there exists $C>0$ such that\vspace*{-3pt}
\begin{equation}\label{eq:estih1}
C\geq \sup_{\substack{k\in\{1,\dots,K\}\\ m\geq m_1}} \norm{ \varphi_{k,m}}_{H^1(\Omega)}.
\end{equation}
Note also that by \eqref{eq:firstest}, \eqref{eq:estih1} and the trace theorem, we get that\vspace*{-3pt}
\begin{equation}\label{eq:estil2tot}
\bigl|\,\norm{ \varphi_{k,m}}_{L^2(\Omega')}-\norm{u_{m+m_0}( \varphi_{k,m})}_{L^2(\Omega')}\bigr|\leq \norm{ \varphi_{k,m} - u_{m+m_0}( \varphi_{k,m})}_{L^2(\Omega')}\leq C/m.
\end{equation}
Moreover, by Proposition \ref{prop:ext}, we obtain that
\[
\norm{u_{m+m_0}( \varphi_{k,m})}_{L^2(\Omega')}^2
\leq (m+m_0)^{-2}\Lambda_{m+m_0}( \varphi_{k,m})\leq C(m+m_0)^{-1}\norm{ \varphi_{k,m}}_{H^1(\Omega)}^2,
\]
and we deduce that\vspace*{-3pt}
\begin{equation}\label{eq:estiextnorm}
\norm{ \varphi_{k,m}}_{L^2(\Omega')}\leq Cm^{-1}.
\end{equation}
Combining \eqref{eq:firstest}, \eqref{eq:estih1}, \eqref{eq:estiextnorm}, and Proposition \ref{lem:asym_eigen_int} with an induction procedure as in the proof of Lemma \ref{lem:induc1}, we get the following result.

\begin{lemma}\label{lem:induc2}
Let $K\in\NN$.
The following properties hold:
\begin{enumerate}
\item\label{pt:ind_tot_asymp1} For all $j\in \{1,2,\dots, K\}$, we have
$\lim_{m\to +\infty} \lambda_{j,m} = \lambda_j$.
\item\label{pt:ind_tot_asymp2} For all subsequence $(m_n)_{n\in \NN}$ going to $+\infty$ as $n\to+\infty$, all $L^2$-orthonormal family of eigenvectors $( \varphi_{1,m_n}, \dots, \varphi_{K,m_n})$ of $|H_m|$ associated with $(\lambda_{1,m_n},\dots, \lambda_{K,m_n})$ such that the sequence $( \varphi_{1,m_n}, \dots, \varphi_{K,m_n})_{n\in \NN}$ converges weakly in $H^1(\Omega)$, we have that the sequence $( \varphi_{1,m_n}, \dots, \varphi_{K,m_n})_{n\in \NN}$ converges strongly in $H^1(\Omega)$ and
\begin{equation}\label{eq.lim_tot_bnd}
\lim_{n\to+\infty}m_{n}\norm{\Xi^- \varphi_{j,m_n}}_{L^2(\Gamma)}^2 = 0
\end{equation}
for all $j\in\{1,\dots,K\}$.
\item\label{pt:ind_tot_asymp3} Every weak limit $(\varphi_1,\dots,\varphi_K)$ of such a sequence is an $L^2$-orthonormal family of eigenvectors of $|H^\Omega|$ associated with the eigenvalues $(\lambda_1,\dots,\lambda_K)$.
\end{enumerate}
\end{lemma}

\begin{remark}\label{rem.convergencefunctions}
In other words, Lemma \ref{lem:induc2} shows the convergence of the eigenspaces associated with the $K$ first eigenvalues of $|H_m|$. Indeed, for all converging subsequence, the corresponding eigenprojector converges to the eigenprojector of $|H^\Omega|$. Thus, when $m$ goes to $+\infty$, the eigenprojector associated with the $K$ first eigenvalues of $|H_m|$ converges to the one of $|H^\Omega|$ associated to the $K$ first eigenvalues. Of course, we have no such convergence result for the individual eigenfunctions.
\end{remark}

\subsection{Second term in the asymptotic}
In this section, we will freely use the following regularity result, whose proof is given in the appendix.

\begin{lemma}\label{lem:regu}
There exists a constant $C>0$ such that for every $m\in \RR$ and every eigenfunction $u$ of $H_m$ associated with an eigenvalue $\lambda\in \RR$, we have
\[
\|u\|_{H^2(\Omega)}\leq C(1+|\lambda|)\|u\|_{L^2(\RR^3)}.
\]
Moreover, for every eigenfunction $u$ \resp $v$ of $H^\Omega$ \resp $L^\intt_m$ associated with an eigenvalue $\lambda\in \RR$, \resp $\lambda^2\in \RR$, we also have that
\begin{align*}
\|u\|_{H^2(\Omega)}&\leq C(1+|\lambda|)\|u\|_{L^2(\Omega)}\\
\tag*{and}
\|v\|_{H^2(\Omega)}&\leq C(1+|\lambda|)\|v\|_{L^2(\Omega)}.
\end{align*}
\end{lemma}

\subsubsection{Upper bound}
In this section, we prove the following lemma.
\begin{lemma}
Let $\lambda$ be an eigenvalue of $|H^\Omega|$ of multiplicity $k_1\in \NN^*$. Let $k_0\in \NN$ be the unique integer such that
\[
\lambda = \lambda_{k_0+1} = \dots = \lambda_{k_0+k_1}.
\]
Then
\begin{equation}\label{eq:bsconclusion}
\limsup_{m\to+\infty} m(\lambda_{k_0+k,m}^2-\lambda^2)\leq\widetilde\nu_{\lambda,k},
\end{equation}
where, for $k\in \{1,\dots,k_1\}$,
\begin{equation}\label{eq:bsminmax}
\widetilde \nu_{\lambda,k}
:=
\inf_{\substack{
V\subset \ker(|H^\Omega|-\lambda \Id)\\
\dim V = k}}\;
\sup_{\substack{v\in V\\
\norm{v}_{L^2(\Omega)}=1}}
\widetilde \eta_\lambda(v)
\end{equation}
and
\[
\widetilde \eta_\lambda(v) := \int_\Gamma \Bigl(\frac{|\nabla_s v|^2}{2} -\frac{|(\pa_\bn+\kappa/2+m_0)v|^2}{2}+ \Bigl(\frac{K}{2} - \frac{\kappa^2}{8}-\frac{\lambda^2}{2}\Bigr)|v|^2\Bigr)\dx \Gamma.
\]
\end{lemma}

\begin{proof}
Let $(u_{1,m}, \dots,u_{k_0+k_1,m})$ be an $L^2$-orthonormal family of eigenvectors of $L^\intt_m$ associated with the eigenvalues $(\lambda^\intt_{1,m},\dots,\lambda^\intt_{k_0+k_1,m})$.
Let $(m_n)_{n\in \NN}$ be a subsequence which goes to $+\infty$ as $n$ tends to $+\infty$ and which satisfies
\begin{enumerate}
\item $\limsup_{m\to+\infty}m(\lambda_{k_0+k,m}^2-\lambda^2) = \lim_{n\to+\infty}m_n(\lambda_{k_0+k,m_n}^2-\lambda^2)$,
\item $(u_{1,m_n},\dots,u_{k_0+k_1,m_n})$ converges in $L^2(\Omega)$ to $(u_1,\dots,u_{k_0+k_1})$,
\end{enumerate}
where $(u_1,\dots,u_{k_0+k_1})$ is an $L^2$-orthonormal family of eigenvectors of $H^\Omega$ associated with the eigenvalues $(\lambda_1,\dots,\lambda_{k_0+k_1})$.
By Lemma \ref{lem:regu}, this sequence is uniformly bounded in $H^2(\Omega)$. By interpolation, the convergence also holds in $H^s(\Omega)$ for all $s\in[0,2)$.

Since \eqref{eq:bsminmax} is a finite dimensional spectral problem, there exists an $L^2$-orthonormal basis $(w_{k_0+1},\dots,w_{k_0+k_1})$ of $\ker(|H^\Omega|-\lambda\Id)$ such that
\[
\widetilde{\eta}_\lambda\biggl(\sum_{s=k_0+1}^{k_0+k_1}a_sw_s\biggr) = \sum_{s=k_0+1}^{k_0+k_1}|a_s|^2\widetilde{\eta}_\lambda(w_s) = \sum_{s=k_0+1}^{k_0+k_1}|a_s|^2\widetilde \nu_{\lambda,s-k_0},
\]
for all $a_{k_0+1},\dots,a_{k_0+k_1}\in \CC$.
Moreover, we have
\[
\ker(|H^\Omega|-\lambda\Id) = \Span(u_{k_0+1},\dots u_{k_0+k_1}) = \Span(w_{k_0+1},\dots w_{k_0+k_1}),
\]
so that there exists a unitary matrix $B\in \CC^{k_1\times k_1}$ such that $Bu = w$, where $u = (u_{k_0+1},\dots,u_{k_0+k_1})^T$ and $w = (w_{k_0+1},\dots,w_{k_0+k_1})^T$.
Using Proposition \ref{prop:ext}, we extend these functions outside $\Omega$ by
\[
\widetilde u_{j,m} = \begin{cases}
u_{j,m}& \mbox{ on } \Omega,\\
u_{m+m_0}(u_{j,m}) &\mbox{ on } \Omega',
\end{cases}
\]
for $j\in\{1,\dots,k_0+k_1\}$.
We also define
\begin{align*}
u_m :&= (u_{k_0+1,m},\dots,u_{k_0+k_1,m})^T,\\
w_{m} &= (w_{k_0+1,m},\dots,w_{k_0+k_1,m})^T := Bu_m,\\
\widetilde w_{m} &= (\widetilde w_{k_0+1,m},\dots,\widetilde w_{k_0+k_1,m})^T := B(\widetilde u_{k_0+1,m},\dots,\widetilde u_{k_0+k_1,m})^T,\\
\tag*{and}
V_{k_0+k,m} &= \Span(u_{1,m},\dots u_{k_0,m}, w_{k_0+1,m},\dots, w_{k_0+k}),
\\
\widetilde{V}_{k_0+k,m} &= \Span( \widetilde{u}_{1,m},\dots \widetilde{u}_{k_0,m}, \widetilde{w}_{k_0+1,m},\dots, \widetilde{w}_{k_0+k}),
\end{align*}
for all $k\in\{1,\dots,k_1\}$ and all $m\geq m_1$. Let us remark that
\[
\dim V_{k_0+k,m} = \dim \widetilde{V}_{k_0+k,m} = k_0+k
\]
for all $k\in\{1,\dots,k_1\}$ (choosing if necessary a larger constant $m_1>0$).
In the following, we consider test functions of the form
\[
v_m = \sum_{j=1}^{k_0}a_j\widetilde u_{j,m}+\sum_{j=k_0+1}^{k_0+k_1}a_j\widetilde w_{j,m},
\]
where $a_{1},\dots, a_{k_0+k_1}\in\CC$ satisfy $\sum_{j=1}^{k_0+k_1}|a_j|^2 = 1$, so that
\[
\|v_m\|_{L^2(\Omega)}^2 = \sum_{j=1}^{k_0+k_1}|a_j|^2 = 1.
\]
By Proposition \ref{prop:ext}, we have
\begin{equation}\label{eq:bs2}
\norm{v_m}_{L^2(\RR^3)}^2 = \norm{v_m}_{L^2(\Omega)}^2 + \norm{v_m}_{L^2(\Omega')}^2
= 1 + \frac{\norm{v_m}_{L^2(\Gamma)}^2}{2m} + \scrO(m^{-2}),
\end{equation}
and
\begin{multline}\label{eq:bs3}
\norm{H_m v_m}_{L^2(\RR^3)}^2 \\
= \cQ^\intt_{m}(v_m) + m^{-1}\int_\Gamma \Bigl(\frac{|\nabla_s v_m|^2}{2} + \Bigl(\frac{K}{2} - \frac{\kappa^2}{8}\Bigr)|v_m|^2\Bigr)\dx \Gamma + \scrO(m^{-3/2}).
\end{multline}
From \eqref{eq:bs2} and \eqref{eq:bs3}, we deduce that
\begin{multline*}
m\Bigl(\frac{\norm{H_m v_m}_{L^2(\RR^3)}^2}{\norm{v_m}_{L^2(\RR^3)}^2}-\lambda^2\Bigr)
\leq m\left(\cQ^\intt_m(v_m)-\lambda^2\right)
\\
+ \int_\Gamma \Bigl(\frac{|\nabla_s v_m|^2}{2} + \Bigl(\frac{K}{2} - \frac{\kappa^2}{8}- \frac{\cQ_m^\intt(v_m)}{2}\Bigr)|v_m|^2\Bigr)\dx \Gamma + \scrO(m^{-1/2}).
\end{multline*}
Then, for $k\in\{1,\dots,k_1\}$, we get
\begin{equation}\label{eq:bs1}\begin{split}
m\big(\lambda_{k_0+k,m}^2-\lambda^2\big)&
\leq
\sup_{v_m\in \widetilde{V}_{k_0+k,m}\setminus \{0\}}
m\Bigl(\frac{\norm{H_m v_m}_{L^2(\RR^3)}^2}{\norm{v_m}_{L^2(\RR^3)}^2}-\lambda^2\Bigr)
\\&
\leq
\sup_{\substack{v_m\in {V}_{k_0+k,m}\\
\norm{v_m}_{L^2(\Omega)}=1}}
m\left(\cQ^\intt_m(v_m)-\lambda^2\right) + \overline{\eta}_m(v_m) + \scrO(m^{-1/2}),
\end{split}\end{equation}
where
\[
\overline{\eta}_m(v) := \int_\Gamma \Bigl(\frac{|\nabla_s v|^2}{2} + \Bigl(\frac{K}{2} - \frac{\kappa^2}{8}-\frac{\cQ_m^\intt(v)}{2}\Bigr)|v|^2\Bigr)\dx \Gamma.
\]

The remaining of the proof concerns the asymptotic behavior of
\[
\mu_{k,m} := \sup_{\substack{v_m\in {V}_{k_0+k,m}\\
\norm{v_m}_{L^2(\Omega)}=1}}
m\left(\cQ^\intt_m(v_m)-\lambda^2\right) + \overline{\eta}_m(v_m),
\]
for $k\in\{1,\dots, k_1\}$ when $m$ goes to $+\infty$. Let us first remark that for every $v_m\in V_{k_0+k,m}$, we have
\[\begin{split}
&v_m = \sum_{j=1}^{k_0}a_j u_{j,m}+\sum_{j=k_0+1}^{k_0+k}a_j w_{j,m} = \sum_{j=1}^{k_0}a_j u_{j,m}+\sum_{s=k_0+1}^{k_0+k_1}\biggl(\sum_{j=k_0+1}^{k_0+k}a_jb_{j,s} \biggr)u_{s,m},
\end{split}\]
where $(b_{j,s})_{j,s\in \{k_0+1,\dots,k_0+k_1\}} = B$.
Thanks to Proposition \ref{lem:asym_eigen_int}, we obtain
\begin{multline}\label{eq:bs4}
m_n\big(\cQ^\intt_{m_n}(v_{m_n})-\lambda^2\big)
\\
\hspace*{3mm}= \sum_{j=1}^{k_0}m_n(\lambda_{j,m_n}^\intt-\lambda^2)|a_j|^2
+
\sum_{j=k_0+1}^{k_0+k_1}m_n(\lambda_{j,m_n}^\intt-\lambda^2)\biggl|\sum_{s=k_0+1}^{k_0+k}a_sb_{s,j} \biggr|^2
\\
=\sum_{j=1}^{k_0}m_n(\lambda_{j,m_n}^\intt\!\!-\!\lambda^2)|a_j|^2
-\frac{\bigl\|\left(\pa_n+\kappa/2+m_0\right)\sum_{j=k_0+1}^{k_0+k}a_j w_{j}\bigr\|_{L^2(\Gamma)}^2}{2}+o(1).
\end{multline}
Using \eqref{eq:bs1} and \eqref{eq:bs4}, and taking $a_1= \dots= a_{k_0+k-1} = 0$, $a_{k_0+k}= 1$, we deduce that
\begin{equation}\label{eq:bs5}
\liminf_{n\to+\infty}\mu_{k,m_n}\geq \widetilde \nu_{\lambda,k}.
\end{equation}
Let $(v^n)_{n\in \NN}$ be a maximizing sequence of $\mu_{k,m_n}$. For all $n$, there exists a unitary vector $a^n =(a_{1,n},\dots, a_{k_0+k,n})\in \CC^{k_0+k} $ such that
\[
v^n = \sum_{j=1}^{k_0}a_{j,n}u_{j,m_n} + \sum_{j=k_0+1}^{k_0+k}a_{j,n}w_{j,m_n}.
\]
Up to a subsequence, we can assume that $(a^n)$ converges in $\CC^{k_0+k}$ to a unitary vector $a = (a_{k_0+1},\dots,a_{k_0+k})$.
Then, Proposition \ref{lem:asym_eigen_int}, \eqref{eq:bs4}, and \eqref{eq:bs5} ensure that
\[
\lim_{n\to+\infty}\lambda_{j,m_n}^\intt-\lambda^2 \leq \lambda_j^2-\lambda^2<0
\] for $j\in\{1,\dots,k_0\}$, thus there exists $c_0>0$ such that
\[
m_n\sum_{j=1}^{k_0}|a_{j,n}|^2\leq c_0\,
\]
and
\[
\limsup_{n\to+\infty}\mu_{k,m_n}\leq \widetilde \eta_\lambda(v)\leq \widetilde \nu_{\lambda,k},
\]
where $v = \sum_{j=k_0+1}^{k_0+k}a_jw_j$. Thanks to \eqref{eq:bs1}, and noticing that $\lim_{n\to+\infty}\mu_{k,m_n}=\widetilde \nu_{\lambda,k}$ and
\[\limsup_{m\to+\infty} m(\lambda_{k_0+k,m}^2-\lambda^2)\leq\widetilde\nu_{\lambda,k},
\]
we conclude the proof.
\end{proof}
\subsubsection{Lower bound}
Let $\lambda$ be the first eigenvalue of $|H^\Omega|$, whose multiplicity is denoted by $k_1\in \NN^*$:
\[
\lambda = \lambda_{1} = \dots = \lambda_{k_1}.
\]

In the following, we look for the second term in the asymptotic expansion of $\lambda$. More precisely, we will show the following result.
\begin{lemma}
For all $k\in \{1,\dots,k_1\}$, we have that
\[
\liminf_{m\to+\infty}m(\lambda_{k,m}^2-\lambda^2)
\geq \widetilde \nu_{\lambda_1,j},
\]
where $\widetilde \nu_{\lambda_1,j}$ is defined in \eqref{eq:bsminmax}.
\end{lemma}
\begin{proof}
By Lemma \ref{lem:induc2} and Proposition \ref{lem:asym_eigen_int}, we have
\[
\lim_{m\to+\infty}\lambda_{k,m}^2 = \lim_{m\to+\infty}\lambda^\intt_{k,m} = \lambda^2,
\]
for all $k\in \{1,\dots,k_1\}$.
Let $(\varphi_{1,m},\dots,\varphi_{k_1,m})$ be an $L^2$-orthonormal family of eigenvectors of $|H_m|$ associated with the eigenvalues $(\lambda_{1,m},\dots,\lambda_{k_1,m})$ for all $m\geq m_1$.
By~Lemma \ref{lem:regu}, there exists $C>0$ such that
\begin{equation}\label{eq:unifboundH2int}
C\geq \sup_{\substack{m\geq m_1
\\
j\in\{1,\dots,k_1\}}}
\norm{\varphi_{j,m}}_{H^2(\Omega)}.
\end{equation}
We remark that, for all $k\in \{1,\dots,k_1\}$, and all $m\geq m_1$,
\[\begin{split}
&
\lambda_{k,m}^2
=
\norm{H_m \varphi_{k,m}}_{L^2(\RR^3)}^2
=
\sup_{\substack{(a_{1},\dots,a_{k})\in\CC^{k}
\\
\sum_{j = 1}^{k}|a_j|^2 = 1}}
\biggl\|H_m \biggl(\sum_{j=1}^{k}a_j\varphi_{j,m}\biggr)\biggr\|_{L^2(\RR^3)}^2.
\end{split}\]
Let $a = (a_{1},\dots,a_{k})\in\CC^{k}$ be such that $\sum_{j = 1}^{k}|a_j|^2 = 1$. We define
\[
\varphi_{m}^a = \sum_{j=1}^{k}a_j\varphi_{j,m}.
\]
Combining \eqref{eq:quadIM}, \eqref{eq:unifboundH2int}, and Proposition \ref{prop:ext}, we get
\begin{multline}\label{eq:secest}
\lambda_{k,m}^2
\geq \cQ^\intt_m\left(\varphi_{m}^a\right)
+ m^{-1}\int_\Gamma \Bigl(\frac{|\nabla_s \varphi_{m}^a|^2}{2} + \Bigl(\frac{K}{2} - \frac{\kappa^2}{8}\Bigr)|\varphi_{m}^a|^2\Bigr)\dx \Gamma
\\
+(m+m_0)^2\norm{\varphi_{m}^a - u_{m+m_0}(\varphi_{m}^a)}_{L^2(\Omega')}^2
+\scrO(m^{-3/2}).
\end{multline}
By \eqref{eq:estil2tot}, we have that
\[\begin{split}
\big|\norm{\varphi_{m}^a}_{L^2(\Omega')}^2&-\norm{u_{m+m_0}(\varphi_{m}^a)}_{L^2(\Omega')}^2\bigr|
\\&\qquad
\leq
C/m\big(\norm{\varphi_{m}^a}_{L^2(\Omega')}+\norm{u_{m+m_0}(\varphi_{m}^a)}_{L^2(\Omega')}\bigr)
\\
&\qquad
\leq C/m\big(\norm{\varphi_{m}^a - u_{m+m_0}(\varphi_{m}^a)}_{L^2(\Omega')}+2\norm{u_{m+m_0}( \varphi_{m}^a)}_{L^2(\Omega')}\bigr)
\\
&\qquad
\leq C/m\bigl(m^{-1}+2\norm{u_{m+m_0}( \varphi_{m}^a)}_{L^2(\Omega')}\bigr).
\end{split}\]
In addition, using Proposition \ref{prop:ext} and \eqref{eq:unifboundH2int}, we deduce that
\[
\Big|
\norm{u_{m+m_0}(\varphi_{m}^a)}_{L^2(\Omega')}^2 - \frac{\norm{\varphi_{m}^a}_{L^2(\Gamma)}^2}{2m}
\Big|
\leq \frac{C}{m^{3/2}}.
\]
Therefore,
\begin{equation}\label{eq:normL2ext}
\Big|\norm{\varphi_{m}^a}_{L^2(\Omega')}^2- \frac{\norm{\varphi_{m}^a}_{L^2(\Gamma)}^2}{2m}\Big|\leq \frac {C}{m^{3/2}}.
\end{equation}
Thanks to \eqref{eq:secest} and Proposition \ref{lem:asym_eigen_int}, we obtain
\begin{multline}\label{eq:thiest}
m(\lambda_{k,m}^2-\lambda^2)
\geq m\left(\cQ^\intt_m(\varphi_{m}^a)-\lambda^2\norm{\varphi_{m}^a}_{L^2(\Omega)}^2\right)
\\
+ \int_\Gamma \Bigl(\frac{|\nabla_s\varphi_{m}^a|^2}{2} + \Bigl(\frac{K}{2} - \frac{\kappa^2}{8}-\frac{\lambda^2}{2}\Bigr)|\varphi_{m}^a|^2\Bigr)\dx \Gamma
+\scrO(m^{-1/2}).
\end{multline}

Let $(u_{j,m})_{j\in \NN^*}$ be an $L^2$-orthonormal basis of $L^2(\Omega;\CC^4)$ whose elements are eigenvectors of $L^\intt_m$ associated with the sequence of eigenvalues $(\lambda_{j,m}^\intt)$.
Since $\lambda_{j,m}^\intt$ converges to $\lambda_j^2$ as $m$ goes to $+\infty$, we get that
\[
\lambda_{j,m}^\intt-\lambda^2\geq 0
\]
for all $j\geq k_1+1$ and all $m\geq m_1$ (choosing if necessary a larger constant $m_1>0$).
We then deduce that
\begin{equation}\label{eq:limlb2}\begin{split}
m\bigl(\cQ^\intt_m(\varphi_{m}^a)-\lambda^2\norm{\varphi_{m}^a}_{L^2(\Omega)}^2\bigr)
&=
\sum_{s=1}^{+\infty}
m\left(\lambda_{s,m}^\intt-\lambda^2\right)|{\braket{\varphi_m^a, u_{s,m}}_{\Omega}}|^2
\\&
\geq
\sum_{s=1}^{k_1}
m\left(\lambda_{s,m}^\intt-\lambda^2\right)|{\braket{\varphi_m^a, u_{s,m}}_{\Omega}}|^2.
\end{split}\end{equation}

Let $(m_n)_{n\in \NN^*}$ be a subsequence which goes to $+\infty$ as $n$ tends to $+\infty$ and such that
\begin{enumerate}
\item $\liminf_{m\to+\infty}m(\lambda_{k,m}^2-\lambda^2) = \lim_{n\to+\infty}m_n(\lambda_{k,m_n}^2-\lambda^2)$,
\item $(u_{1,m_n},\dots,u_{k_1,m_n})$ converges in $H^1(\Omega)$ to $(u_1,\dots,u_{k_1})$,
\item $(\varphi_{1,m_n},\dots,\varphi_{k_1,m_n})$ converges in $H^1(\Omega)$ to $(\varphi_1,\dots,\varphi_{k_1})$, \end{enumerate}
where $(u_1,\dots,u_{k_1})$ and $(\varphi_1,\dots,\varphi_{k_1})$ are $L^2$-orthonormal families of eigenvectors of~$H^\Omega$ associated with the eigenvalue $\lambda$.
By Proposition \ref{lem:asym_eigen_int}, we have that
\begin{equation}\label{eq:limlbfin}\begin{split}
\lim_{n\to+\infty}\sum_{s=1}^{k_1}
m\left(\lambda_{s,m_n}^\intt-\lambda^2\right)&|{\braket{\varphi_{m_n}^a, u_{s,m_n}}_{\Omega}}|^2
\\&
=
\sum_{s=1}^{k_1}
-\frac{\norm{(\pa_\bn+\kappa/2+m_0)u_s}_{L^2(\Gamma)}^2}{2}\,|{\braket{\varphi^a, u_{s}}_{\Omega}}|^2\\
&=
-\frac{\norm{(\pa_\bn+\kappa/2+m_0)\varphi_a}_{L^2(\Gamma)}^2}{2},
\end{split}\end{equation}
where $\varphi^a = \sum_{j=1}^{k}a_j\varphi_j$.
From \eqref{eq:thiest}, \eqref{eq:limlb2}, and \eqref{eq:limlbfin}, we obtain
\[
\liminf_{m\to+\infty}m(\lambda_{k,m}^2-\lambda^2)
\geq \widetilde \eta_\lambda(\varphi^a)\,
\]
and
\[
\liminf_{m\to+\infty}m(\lambda_{k,m}^2-\lambda^2)
\geq
\sup_{\substack{(a_{1},\dots,a_k)\in \CC^{k}
\\
\sum_{j=1}^{k}|a_j|^2 = 1}}
\widetilde \eta_\lambda(\varphi^a)
\geq \widetilde\nu_{\lambda,k}.
\]
Then, the conclusion follows from this and the upper bound \eqref{eq:bsconclusion}.
\end{proof}
\begin{remark}
When considering a larger eigenvalue $\lambda > \lambda_{1}$, the proof above breaks down since
\[
\sum_{s=1}^{k_0}
m\left(\lambda_{s,m}^\intt-\lambda^2\right)|{\braket{\varphi_m^a, u_{s,m}}_{\Omega}}|^2
\]
is non positive and the non-wanted terms in \eqref{eq:limlb2} cannot be removed so easily anymore. In the expression above, $k_0$ denotes the unique integer such that
\[
\lambda = \lambda_{k_0+1} = \dots = \lambda_{k_0+k_1}.
\]
\end{remark}

\appendix
\section*{Appendix. Sketch of the proof of Lemma \ref{lem:regu}}\label{sketch lm}
\refstepcounter{section}

The purpose of this appendix is to give the main ideas of the proof of Lemma \ref{lem:regu}. We recall that the boundary is supposed to have $\cC^2$ regularity. We do not intend to give a rigorous proof but rather to enlighten why the classical arguments give uniform bounds in $m$ (see for instance \cite[\S 6.3]{evans1998partial}). In particular, we restrict ourselves to the operator $H_m$ for $\Omega = \RR^3_+ := \{\bx = (x_1,x_2,x_3) \sep x_3>0\}$, and we consider the solution $u\in H^1(\RR^3; \CC^4)$ of
\[
H_m u = (\alpha\cdot D +(m_0+m\chi_{\RR^3_-})\beta)u = f,
\]
where $f\in H^1(\RR^3;\CC^4)$.
By Lemma \ref{lem:quadf1} and Proposition \ref{lem:ext}, we have
\begin{multline*}
\norm{f}_{L^2(\RR^3)}^2\geq \norm{\nabla u}_{L^2(\Omega)}^2 +m_0^2\norm{u}_{L^2(\Omega)}^2 +m_0\norm{u}_{L^2(\Gamma)}^2+ \sum_{k=1}^2\norm{\pa_k u}_{L^2(\Omega')}^2 \\
+ 2m\norm{\Xi^- u}_{L^2(\Gamma)}^2 - C/m\norm{u}_{L^2(\Gamma)}^2,
\end{multline*}
so that by the trace theorem, there exists $C>0$ such that
\begin{equation}\label{lem.reguesti1}\begin{split}
&C\big(\norm{f}_{L^2(\RR^3)}^2+\norm{u}_{L^2(\Omega)}^2\bigr)\geq \norm{\nabla u}_{L^2(\Omega)}^2 + \sum_{k=1}^2\norm{\pa_k u}_{L^2(\Omega')}^2.
\end{split}\end{equation}
Using the notation of \cite[\S 6.3]{evans1998partial}, we introduce the difference quotients
\[
D^h_ku(\bx) = \frac{u(\bx+he_k)-u(\bx)}{h}, \quad h\in \RR,\; h\ne 0,\; \bx\in\RR^3,\; k\in \{1,2,3\}.
\]
For $j\in \{1,2\}$, we get that
\[
H_m D^h_j u = (\alpha\cdot D +(m_0+m\chi_{\RR^3_-})\beta)D^h_j u =D^h_j f,
\]
and then, using \eqref{lem.reguesti1}, we obtain
\begin{equation*}
C\big(\norm{D^h_j f}_{L^2(\RR^3)}^2+\norm{D^h_j u}_{L^2(\Omega)}^2\bigr)\geq \norm{\nabla D^h_ju}_{L^2(\Omega)}^2 + \sum_{k=1}^2\norm{\pa_k D^h_j u}_{L^2(\Omega')}^2.
\end{equation*}
By \cite[\S 5.8.2]{evans1998partial}, we deduce that
\begin{multline}\label{lem.reguesti2}
C\big(\norm{\pa_j f}_{L^2(\RR^3)}^2+\norm{\pa_j u}_{L^2(\Omega)}^2+\norm{ f}_{L^2(\RR^3)}^2+\norm{ u}_{L^2(\Omega)}^2\bigr)\\\geq \norm{\nabla \pa_ju}_{L^2(\Omega)}^2 + \sum_{k=1}^2\norm{\pa_k \pa_j u}_{L^2(\Omega')}^2.
\end{multline}
We also have that, in $\Omega$,
\[
-\pa_3^2u = H_m^2u+\biggl(\sum_{k=1}^2\pa_k^2-m_0^2\biggr)u = H_mf,
\]
thus
\begin{equation}\label{lem.reguesti3}
\norm{\pa_3^2 u}_{L^2(\Omega)}\leq C\norm{f}_{H^1(\Omega)}.
\end{equation}
Using \eqref{lem.reguesti1}, \eqref{lem.reguesti2} and \eqref{lem.reguesti3}, we get the desired estimate.


\bibliographystyle{jepplain}
\bibliography{arrizabalaga-et-al}
\end{document} 
